\documentclass[12pt,letterpaper]{article}

\usepackage{times}
\usepackage{graphicx}
\usepackage{amsmath}
\usepackage{amssymb}
\usepackage{bm}
\usepackage{cite}
\usepackage{hyperref}
\usepackage{microtype}
\usepackage{booktabs}
\usepackage{geometry}
\usepackage{setspace}

\geometry{letterpaper, margin=1in}
\doublespacing

\newcommand{\revision}[1]{#1}


\title{\textbf{Consolidated Report: Effective Trace Framework for Self-Similar Casimir Systems, FDTD Scaling Results, and Semiclassical Gravity Backreaction}}

\author{\textbf{Goren Gordon} \\
Department of Informatics, Indiana University Bloomington \\
Bloomington, Indiana, USA \\
E-mail: \href{mailto:goren@gorengordon.com}{goren@gorengordon.com}}

\date{\today}

\begin{document}

\maketitle

\begin{abstract}
The interaction of quantum fields with fractal and self-similar geometries encompasses multiple distinct physical regimes, including spectral geometry on intrinsic fractals, macroscopic self-similar Casimir configurations, and bounded Euclidean cavities with fractal boundaries.
While the thermal equations of state and spectral asymptotics for these systems are well established, a cohesive treatment of the vacuum trace frequently conflates rigorous mathematical bounds with phenomenological models.
In this manuscript, we systematically decouple these regimes and advance a unified effective framework combining the rigorous thermal trace of fractal radiation with a zero-temperature integrated vacuum trace for plate-like self-similar geometries.
We demonstrate that for systems governed by a scale-dependent Casimir coefficient $C(d_s, \ln(d/\ell_*))$, the anisotropic stress-energy tensor produces an integrated vacuum trace proportional to its logarithmic running, $\partial_{\ln d}C$.
We strictly differentiate this effective macroscopic backreaction from first-principles local trace anomalies on genuine fractal boundaries.
To bridge the threshold toward experimental verification, we evaluate this framework across multiple operational tiers: first, through macroeconomic thoughts and localized atomic probes (including fractal Purcell, Casimir-Polder, and geometric Lamb shifts); and second, via coherent macroscopic metamaterial realizations like the Synthetic Gravity Dipole, establishing the analytical prerequisites necessary to transition this effective formalism into a quantitatively predictive electromagnetic theory amenable to physical verification.
\end{abstract}

\newpage
\section{Introduction}
The Casimir effect serves as the quintessential macroscopic manifestation of boundary-dependent quantum vacuum fluctuations \cite{casimir1948}.
Concurrently, extensive developments in analysis on fractals and spectral geometry have demonstrated that self-similar structures fundamentally alter heat kernels, zeta functions, diffusion laws, and thermodynamic equations of state in a manner dictated by the spectral dimension rather than the topological dimension \cite{akkermans2010,kigami2001,kajino2010}.
More recently, self-similar Casimir configurations---including Cantor-like stacks and compartmentalized Sierpi\'nski geometries---have been shown to admit nontrivial recursion relations for vacuum energy, with structural sign changes emerging in select exactly solvable scalar models \cite{shajesh2016,shajesh2017}.

These theoretical developments precipitate a fundamental question: is it possible to rigorously formulate a ``fractal Casimir effect'' that systematically delineates established mathematical results from phenomenological hypotheses?
The resolution is nontrivial, as the literature frequently conflates three distinct physical and geometric configurations: (i) quantum fields propagating on genuine fractal Laplacians;
(ii) Euclidean cavities bounded by fractal boundaries or exhibiting fractal porosity; and (iii) macroscopic self-similar Casimir architectures constructed from standard boundaries.
Each configuration possesses distinct spectral properties and an independent status within theoretical literature.

The scope of this manuscript is defined by four specific objectives: (i) to synthesize the established thermal, spectral, and Casimir formalisms applicable to fractal systems;
(ii) to enforce a strict demarcation between mathematically rigorous thermal traces and effective zero-temperature vacuum traces;
(iii) to derive an anisotropic integrated vacuum trace for plate-like self-similar geometries governed by a logarithmically running Casimir coefficient; and
(iv) to explicitly transition these theoretical conjectures into precision experimental and observational metrics through microscopic atomic sensors and coherent macroscopic geometries.

The manuscript is structured as follows. Section~\ref{sec:prior} systematically reviews the theoretical background. Section~\ref{sec:novel} articulates the proposed effective framework, exploring both its thermal and vacuum formulations alongside finite-temperature generalizations. Section~\ref{sec:gedanken} presents two macroscopic Gedankenexperiments isolating the violation of scale invariance. Section~\ref{sec:atomic_probes} maps the framework onto microscopic test systems, evaluating modified atomic electrodynamics and weak semiclassical gravitational signatures. Section~\ref{sec:experimental} outlines physical pathways for engineering realization, and Section~\ref{sec:discussion} contextualizes these boundaries before our concluding remarks.

\section{Theoretical Background}
\label{sec:prior}

\subsection{Spectral Geometry and Thermodynamics on Fractals}

In the context of a genuine fractal Laplacian, the small-$\tau$ heat-kernel trace exhibits the asymptotic scaling
\begin{equation}
K_F(\tau) \sim \left(\frac{L_s^2}{\tau}\right)^{d_s/2} F_{\rm per}\!\left(\ln\frac{\tau}{L_s^2}\right),
\label{eq:fractal_heat_kernel}
\end{equation}
where $L_s$ denotes a spectral length scale, $d_s$ represents the spectral dimension, and $F_{\rm per}$ is a multiplicatively periodic modulation function \cite{akkermans2010,kajino2010}.
This log-periodic oscillatory structure is the direct spectral signature of discrete scale invariance.

Akkermans, Dunne, and Teplyaev demonstrated that a massless field confined to such a fractal spatial manifold obeys the thermal equation of state
\begin{equation}
p\,V_s = \frac{U}{d_s},
\label{eq:fractal_eos_prior}
\end{equation}
where $V_s$ is the spectral volume, derived from the residue of the spectral zeta function at $s=d_s/2$ \cite{akkermans2010}.
Within this exact framework, the zero-temperature vacuum energy admits a natural zeta-function regularization:
\begin{equation}
E_0 = \frac{\hbar c}{2L_s}\,\zeta_F\!\left(-\frac{1}{2}\right).
\label{eq:fractal_vacuum_prior}
\end{equation}
We treat Eqs.~\eqref{eq:fractal_eos_prior} and \eqref{eq:fractal_vacuum_prior} as verified foundational constraints.

The operational distinction between Hausdorff and spectral dimensions is paramount.
The Hausdorff dimension $d_h$ parametrizes geometric filling, whereas the spectral dimension governs both diffusion and the asymptotic density of states.
The canonical relation for diffusion on fractals is given by \cite{alexanderorbach1982}
\begin{equation}
d_s = \frac{2d_h}{d_w},
\label{eq:ds_dw}
\end{equation}
where $d_w$ is the walk dimension.
While exact for specific mathematical models, this relation is effectively approximate in physical implementations and should not be treated as a universal identity for arbitrary prefractal resonators.

\subsection{Casimir Energy in Standard and Self-Similar Geometries}

For two ideally conducting parallel plates in $3+1$ dimensions, the Casimir interaction energy per unit area and the associated pressure are given by
\begin{equation}
\frac{E_{\rm flat}(d)}{A_\perp} = -\frac{\pi^2\hbar c}{720\,d^3}, \qquad P_{\rm flat}(d) = -\frac{\pi^2\hbar c}{240\,d^4}.
\label{eq:flat_casimir_prior}
\end{equation}
Brown and Maclay originally derived the renormalized electromagnetic stress tensor for this configuration, demonstrating that the tensor is exactly traceless between the plates in the ideal conformal limit \cite{brownmaclay1969}.
This provides the canonical benchmark against which any hypothesized vacuum trace must be evaluated.

Self-similar Casimir architectures exhibit fundamentally divergent behavior. Shajesh \emph{et al.} formulated several fractal and self-similar scalar-field configurations---notably including a Cantor-set stack of $\delta$-function plates---and demonstrated that the resultant interaction energy can manifest as positive, negative, or zero, strictly dependent on the geometric arrangement \cite{shajesh2016}.
In instances of positive interaction energy, the vacuum fluctuations induce structural expansion.
Furthermore, exact scalar Casimir energies for Sierpi\'nski triangles and rectangles have been derived by exploiting principles of self-similarity and compartmentalization \cite{shajesh2017}.

\revision{To clarify how zero-temperature spatial recursion relations relate to the thermal equation of state, we must examine the joint scaling behavior of the partition function. At zero temperature, spatial recursion relations dictate how the vacuum energy scales with geometry, such as the Cantor-like scaling of parallel plates where the spatial separation $d$ is rescaled, $\Delta e(d) = \Delta e(d/2) + \Delta e_{12}(d)$ \cite{shajesh2016}. However, a purely spatial recursion relation at $T=0$ does not automatically yield the finite-temperature state. To bridge this conceptual gap, we must utilize the Euclidean path-integral partition function $Z(d, \beta)$ defined over periodic imaginary time with period $\beta = \hbar / (k_B T)$ \cite{bordag2009}. In a scale-invariant field theory, a simultaneous dilatation of the spatial boundary $d \to d/\lambda$ and the thermal parameter $\beta \to \beta/\lambda$ leaves the ratio $d/\beta = d T$ invariant. The interaction free energy per unit area $\Delta f(d, T)$ then satisfies the bivariate scaling relation \cite{bordag2009}:
\begin{equation}
\Delta f(d, T) = \lambda^3 \Delta f\left(\frac{d}{\lambda}, \lambda T\right).
\label{eq:bivariate_scaling}
\end{equation}

By taking the zero-temperature limit ($\lambda \to 0$), the effective thermal parameter $\lambda T$ vanishes, the thermal scale decouples, and this joint scaling relation reduces to the standard spatial recursion relation. At finite temperature, the underlying scale invariance of the boundary configuration forces the free energy to scale with the effective dimension of the system (namely, the spectral dimension $d_s$), generating the thermal equation of state $p V_s = U / d_s$ \cite{akkermans2010}. Thus, the thermal state is not directly generated by the spatial recursion itself; rather, both are individual projections of a unified bivariate scaling of the Euclidean partition function under joint spatial and temporal dilations.}

While these precedents establish the possibility of sign-changing Casimir energies in self-similar systems, they must be contextualized against broader constraints.
Kenneth and Klich rigorously proved that any two bodies related by reflection will mutually attract, assuming ordinary conductors or nonmagnetic dielectrics in vacuum \cite{kennethklich2006}.
Therefore, hypotheses suggesting ubiquitous repulsion between arbitrary ``fractal plates'' are physically invalid;
such sign changes are mathematically restricted to non-mirror configurations, intricate many-body geometries, or highly specific scalar models.

\subsection{Fractal Drums, Density of States, and Operator Class}

A critical physical distinction separates intrinsic fractal Laplacians from Euclidean domains bounded by fractal geometries.
For a bounded Euclidean domain $\Omega\subset\mathbb R^n$ possessing a fractal boundary of Minkowski dimension $D$, Lapidus demonstrated that the leading Weyl term retains its standard Euclidean scaling.
The boundary fractality manifests strictly through the remainder estimate of the spectral asymptotics \cite{lapidus1991}.
Schematically, this is expressed as
\begin{equation}
N_\Omega(\lambda)=C_n|\Omega|\lambda^{n/2}+\mathcal O(\lambda^{D/2}),
\label{eq:lapidus_schematic}
\end{equation}
subject to the conditions formulated in Ref.~\cite{lapidus1991}.
Thus, mapping a field onto a Koch snowflake domain does not override the leading Euclidean exponent with a non-integer spectral exponent.

Conversely, intrinsic fractal Laplacians unconditionally exhibit the spectral scaling
\begin{equation}
N(\lambda) \sim \lambda^{d_s/2}G(\ln\lambda),
\label{eq:genuine_fractal_counting}
\end{equation}
modulated by a multiplicatively periodic function $G$ \cite{akkermans2010,kajino2010}.
Consequently, anomalous exponent scaling is an intrinsic property of the underlying operator class, not merely a byproduct of boundary complexity.
This rigorous distinction underpins the methodology of this paper: the proposed framework targets \emph{plate-like self-similar Casimir geometries and prefractal implementations}, rather than generalized fractal drums.

\subsection{Conformal Parameter Ambiguity and Trace of the Stress-Energy Tensor}

The emergence of a non-vanishing renormalized trace is a well-documented phenomenon.
Within curved-space quantum field theory, the trace anomaly of classically conformal theories is a canonical result \cite{browncassidy1977,birrelldavies1982}.
Herzog and Huang further elucidated the correspondence between vacuum stress tensors and trace anomalies within conformal field theories on curved backgrounds \cite{herzoghuang2013}.

\revision{When formulating the stress-energy tensor for a scalar field, one must confront a fundamental ambiguity associated with the choice of the conformal coupling parameter $\xi$ (often denoted as $\gamma$ in other contexts) \cite{milton2016}. For a real scalar field $\phi$ in a spatially varying background potential $V(z)$, the classical stress-energy tensor is defined by \cite{milton2016}:
\begin{align}
T_{\mu\nu} &= \partial_\mu \phi \partial_\nu \phi - \frac{1}{2} g_{\mu\nu} \partial_\lambda \phi \partial^\lambda \phi \nonumber\\
&\quad - \xi \left( \nabla_\mu \nabla_\nu - g_{\mu\nu} \Box \right) \phi^2 + g_{\mu\nu} \left[ 2\xi V(z) - \frac{1}{2} V(z) \right] \phi^2.
\label{eq:conformal_tensor_classical}
\end{align}
Classically, this tensor is conserved ($\nabla^\mu T_{\mu\nu} = 0$), but its trace is given by \cite{milton2016}:
\begin{equation}
T^\mu{}_\mu = 4 \left( \xi - \frac{1}{6} \right) \Box \phi^2 + \left(2\xi - \frac{1}{2}\right) V(z) \phi^2,
\label{eq:conformal_trace_classical}
\end{equation}
which vanishes in the massless, conformally coupled limit ($\xi = 1/6$) in the absence of a potential. Under ultraviolet regularization, the vacuum expectation value of the stress-energy tensor $\langle T_{\mu\nu} \rangle$ is evaluated by point-splitting in time and transverse directions, introducing a temporal-spatial cutoff $\delta = \sqrt{\tau^2 + \rho^2}$ \cite{milton2016}. Renormalization is performed by WKB subtraction of the divergent terms up to second order \cite{milton2016}:
\begin{equation}
\begin{aligned}
\tilde{g}(z,z; \kappa) \approx & \frac{1}{2\sqrt{\kappa^2 + V(z)}} - \frac{V''(z)}{16(\kappa^2 + V(z))^{5/2}} \\
& + \frac{5 [V'(z)]^2}{64(\kappa^2 + V(z))^{7/2}}.
\end{aligned}
\label{eq:wkb_subtraction}
\end{equation}
This procedure leaves a finite, regularized remainder in the local energy density $\langle T_{00} \rangle_{\rm ren}$ and transverse pressure components that depends explicitly on the arbitrary conformal parameter $\xi$ \cite{milton2016}. This conformal parameter ambiguity must be emphasized, as it suggests that local observables are not uniquely defined. However, the physical normal Casimir pressure and the total integrated Casimir energy remain independent of this parameter. For instance, in a parallel-plate geometry satisfying Dirichlet boundary conditions ($\phi = 0$ on boundaries), the integrated energy density is given by \cite{milton2016,brownmaclay1969}:
\begin{equation}
\int_{0}^{d} \langle T_{00} \rangle_{\rm ren} \, dz = \int_{0}^{d} u_0(z) \, dz + \left( \frac{1}{6} - \xi \right) \left[ \frac{d}{dz} \langle \phi^2 \rangle_{\rm ren} \right]_{0}^{d},
\label{eq:integrated_energy_conformal}
\end{equation}
where $u_0(z)$ is the energy density evaluated at the conformal coupling value $\xi = 1/6$. Under Dirichlet boundary conditions, the field fluctuations are pinned at the boundaries, and the surface term,
\begin{equation}
\left[ \frac{d}{dz} \langle \phi^2 \rangle_{\rm ren} \right]_{0}^{d},
\label{eq:surface_divergence}
\end{equation}
represents a boundary-localized divergence that is independent of the plate separation $d$ or vanishes under appropriate surface-regularization schemes \cite{milton2016}. Because the physical Casimir pressure is obtained via the derivative of the integrated energy with respect to the plate separation $d$, the $\xi$-dependent surface term yields no contribution. Thus, while the local trace $\langle T^\mu{}_\mu \rangle_{\rm ren}$ contains a term proportional to $(1-6\xi)\Box \langle \phi^2 \rangle_{\rm ren}$, the macroscopic, integrated observables of the effective trace framework are completely independent of this conformal parameter ambiguity \cite{milton2016}.}

While the present work adopts this established nomenclature, it does not purport to derive a fundamentally new local anomaly theorem for fractal boundaries.
Instead, the theoretical novelty resides in formulating a specialized integrated vacuum-stress structure governed by scale-dependent Casimir coefficients.

\subsection{Thermodynamic Reconciliation via Mode Tunneling}
A major conceptual challenge in self-similar Casimir systems is the apparent conflict between spatial compartmentalization (which assumes impenetrable Dirichlet boundaries to decouple sub-cavities) and thermodynamics (which assumes uniform thermal equilibrium at temperature $T$). Under our microscopic model of semitransparent barriers modeled by delta potentials of finite coupling $\lambda_0 < \infty$, this paradox is resolved. The transmission coefficient $t(\kappa)$ of an electromagnetic wave tunneling through a single semitransparent plate of potential strength $\lambda_0$ is:
\begin{equation}
|t(\kappa)|^2 = \frac{4\kappa^2}{4\kappa^2 + \lambda_0^2}.
\label{eq:SI_trans}
\end{equation}
For any finite coupling $\lambda_0 < \infty$, the transmission coefficient satisfies $|t(\kappa)|^2 > 0$ for high-frequency modes ($\kappa \gg \lambda_0$). These high-energy modes tunnel through the barriers, establishing a dynamic thermal link between adjacent compartments. This mode tunneling acts as the physical mechanism for heat exchange, allowing the entire partitioned system to thermalize with an external reservoir at a uniform temperature $T$. 

The strict Dirichlet compartmentalization is then recovered as a formal mathematical limit where the potential strength is taken to infinity:
\begin{equation}
\lim_{\lambda_0 \to \infty} |t(\kappa)|^2 = 0.
\label{eq:SI_dirichlet_limit}
\end{equation}
By performing the thermalization step at finite coupling $\lambda_0 < \infty$ and subsequently taking the limit $\lambda_0 \to \infty$, the thermodynamic consistency of self-similar Casimir cavities is rigorously established.

\section{Effective Trace Framework for Self-Similar Geometries}
\label{sec:novel}

\subsection{Scope of the Proposal}

The primary contribution of this work lies not in re-deriving the thermal equation of state \eqref{eq:fractal_eos_prior}, the fractal vacuum energy \eqref{eq:fractal_vacuum_prior}, or the structural sign changes in self-similar Casimir interactions, which are firmly established.
Rather, we introduce a \emph{unified effective trace framework} that coherently integrates: (i) the rigorous thermal trace yielded by fractal thermodynamics and;
(ii) an effective zero-temperature integrated vacuum trace induced by a logarithmically running Casimir coefficient in plate-like self-similar geometries.
This second construct is inherently phenomenological: it is formulated from the integrated energy and anisotropic pressure distributions, distinct from a first-principles derivation of a local renormalized stress tensor on a genuine fractal boundary.

\subsection{Thermal Sector: Exact Trace from Established Fractal Thermodynamics}

To formalize the thermal sector, we define the thermal energy density relative to the spectral volume as
\begin{equation}
\rho_{\rm th}=\frac{U_{\rm th}}{V_s}.
\end{equation}

\revision{To illustrate the explicit dependence of the internal energy on the dimension of space, we derive the thermodynamic variables for a massless scalar field confined to a highly symmetric fractal manifold characterized by a spectral volume $V_s$ and spectral dimension $d_s$ \cite{akkermans2010}. The asymptotic density of states $\mathcal{N}(\omega)$ is governed by the spectral Weyl law on fractals \cite{akkermans2010,kajino2010}:
\begin{equation}
\mathcal{N}(\omega) \approx \frac{V_s}{\Gamma(d_s/2)} \omega^{d_s - 1}.
\label{eq:weyl_spectral_density}
\end{equation}
The thermal partition function $\ln Z(\beta)$ at inverse temperature $\beta = 1/(k_B T)$ is obtained via the standard relation \cite{akkermans2010}:
\begin{equation}
\ln Z(\beta) = -\int_{0}^{\infty} \mathcal{N}(\omega) \ln\left(1 - e^{-\beta \hbar \omega}\right) d\omega.
\label{eq:partition_integral_fractal}
\end{equation}
Expanding the logarithm and performing the integration yields:
\begin{equation}
\ln Z(\beta) = V_s \frac{\Gamma(d_s)}{\Gamma(d_s/2)} \zeta(d_s + 1) \left( \frac{k_B T}{\hbar} \right)^{d_s}.
\label{eq:partition_final_fractal}
\end{equation}
The thermal internal energy $U_{\rm th}(T, d_s)$ is evaluated by differentiating the partition function with respect to $\beta$:
\begin{equation}
U_{\rm th}(T, d_s) = -\frac{\partial}{\partial \beta} \ln Z = d_s V_s \frac{\Gamma(d_s)}{\Gamma(d_s/2)} \zeta(d_s + 1) \frac{(k_B T)^{d_s + 1}}{\hbar^{d_s}},
\label{eq:internal_energy_ds}
\end{equation}
where $\zeta(x)$ is the Riemann zeta function and $\Gamma(y)$ is the Gamma function. Equation~\eqref{eq:internal_energy_ds} explicitly illustrates how the spectral dimension $d_s$ controls both the power-law temperature exponent and the thermodynamic coefficient of the internal energy. In the standard 3D Euclidean limit, setting $d_s = 3$ and replacing the spectral volume $V_s$ with the geometric spatial volume $V$ yields:
\begin{equation}
\text{U}_{\rm th}(T, 3) = 3 V \frac{\Gamma(3)}{\Gamma(3/2)} \zeta(4) \frac{(k_B T)^4}{\hbar^3} = \frac{\pi^2 V (k_B T)^4}{30 (\hbar c)^3},
\label{eq:stefan_boltzmann_limit}
\end{equation}
which recovers the standard Stefan-Boltzmann law \cite{akkermans2010}. On a fractal spatial manifold, the internal energy scales non-trivially as $U_{\rm th} \propto T^{d_s + 1}$.}

Applying the established equation of state \eqref{eq:fractal_eos_prior}, the thermal pressure is determined to be
\begin{equation}
p_{\rm th}=\frac{\rho_{\rm th}}{d_s}.
\label{eq:thermal_pressure_new}
\end{equation}
Adopting the standard mixed-index metric signature, the thermal stress-energy tensor is
\begin{equation}
T^\mu{}_{\nu,\rm th}=\mathrm{diag}(-\rho_{\rm th},p_{\rm th},p_{\rm th},p_{\rm th}),
\end{equation}
yielding the exact trace
\begin{equation}
T^\mu{}_{\mu,\rm th} = -\rho_{\rm th}+3p_{\rm th} = \rho_{\rm th}\left(\frac{3}{d_s}-1\right).
\label{eq:thermal_trace_new}
\end{equation}
Equation~\eqref{eq:thermal_trace_new} represents a mathematically rigorous deduction from established fractal thermodynamics, serving as the first constituent of the unified trace formula.

\subsection{Vacuum Sector: Scale-Dependent Casimir Coefficient and Integrated Trace}

The zero-temperature Casimir regime requires distinct analytical treatment.
For a plate-like configuration, the fundamental observable is the interaction energy per unit transverse area, defined as
\begin{equation}
e(d) \equiv \frac{E_0(d)}{A_\perp}.
\end{equation}
In the smooth Euclidean conformal limit, the energy scales as $e(d)\propto d^{-3}$ with a static coefficient, corresponding to \eqref{eq:flat_casimir_prior}.
Guided by known self-similar Casimir mechanics and the log-periodic modulations characteristic of fractal heat kernels, we propose the generalized ansatz
\begin{equation}
e(d)=\frac{\hbar c}{d^3} C\!\left(d_s,\ln\frac{d}{\ell_*}\right),
\label{eq:edef}
\end{equation}
where $\ell_*$ designates the short-distance crossover scale below which the geometric self-similarity is no longer resolved, and $C$ acts as a dimensionless, scale-dependent effective coefficient.
In the conformal limit, $C$ degenerates to a constant, recovering the standard Casimir scaling.

Applying the principle of virtual work, the normal pressure exerted on the boundary is
\begin{equation}
P_\perp(d)=-e'(d) = \frac{\hbar c}{d^4} \left[ 3C\!\left(d_s,\ln\frac{d}{\ell_*}\right) - \partial_{\ln d}C\!\left(d_s,\ln\frac{d}{\ell_*}\right) \right].
\label{eq:pperp}
\end{equation}
The corresponding vacuum energy density contained within the gap is
\begin{equation}
\rho_{\rm vac}(d)=\frac{e(d)}{d} = \frac{\hbar c}{d^4} C\!\left(d_s,\ln\frac{d}{\ell_*}\right).
\label{eq:rhovac}
\end{equation}
Assuming transverse translational invariance, variation of the transverse area under fixed separation establishes the tangential pressure
\begin{equation}
P_{\parallel}(d)=-\frac{e(d)}{d}=-\rho_{\rm vac}(d).
\label{eq:pparallel}
\end{equation}
Consequently, the vacuum state induces an anisotropic mixed-index stress tensor:
\begin{equation}
T^\mu{}_{\nu,\rm vac} = \mathrm{diag}\!\bigl(-\rho_{\rm vac},P_\parallel,P_\parallel,P_\perp\bigr).
\label{eq:vacuum_tensor_new}
\end{equation}
The trace of this tensor evaluates to
\begin{align}
T^\mu{}_{\mu,\rm vac} &= -\rho_{\rm vac}+2P_\parallel+P_\perp \nonumber\\
&= -\frac{\hbar c}{d^4} \,\partial_{\ln d}C\!\left(d_s,\ln\frac{d}{\ell_*}\right).
\label{eq:vacuum_trace_new}
\end{align}
This result constitutes a central finding. It explicitly demonstrates that the integrated vacuum trace is not dictated by the attractive or repulsive nature of the Casimir force, but rather by the \emph{running} of the effective coefficient across logarithmic scales.
A static $C$ ensures a vanishing trace regardless of force direction;
conversely, a running coefficient fundamentally breaks scale invariance, endowing the vacuum sector with a non-zero integrated trace.

A highly applicable parametrization of the running coefficient is given by
\begin{equation}
C\!\left(d_s,\ln\frac{d}{\ell_*}\right) = C_0(d_s) \left[1+F_{\rm per}\!\left(\ln\frac{d}{\ell_*}\right)\right],
\label{eq:Cdecomp}
\end{equation}
where $F_{\rm per}$ captures the discrete log-periodic scaling.
Under this parametrization, the trace simplifies to
\begin{equation}
T^\mu{}_{\mu,\rm vac} = -\frac{\hbar c\,C_0(d_s)}{d^4} F_{\rm per}'\!\left(\ln\frac{d}{\ell_*}\right).
\label{eq:vacuum_trace_periodic}
\end{equation}
The physical implications are unambiguous: discrete structural self-similarity imposes a residual logarithmic running onto an otherwise scale-invariant vacuum stress, which acts as the explicit generator of the effective integrated trace.

\subsection{Unified Trace and Semiclassical Backreaction}

The total effective trace of the system is the linear superposition of the thermal and vacuum sectors,
\begin{equation}
T^\mu{}_{\mu} = T^\mu{}_{\mu,\rm th}+T^\mu{}_{\mu,\rm vac}.
\label{eq:total_trace_new}
\end{equation}
Substituting \eqref{eq:thermal_trace_new} and \eqref{eq:vacuum_trace_new} yields the unified formula
\begin{equation}
T^\mu{}_{\mu} = \rho_{\rm th}\left(\frac{3}{d_s}-1\right) - \frac{\hbar c}{d^4} \,\partial_{\ln d}C\!\left(d_s,\ln\frac{d}{\ell_*}\right).
\label{eq:total_trace_explicit_new}
\end{equation}
Mapping this trace onto the framework of semiclassical gravity, the traced Einstein field equations,
\begin{equation}
R=-8\pi G\,T^\mu{}_{\mu},
\label{eq:trace_einstein_new}
\end{equation}
dictate the resultant scalar curvature:
\begin{equation}
R = -8\pi G \left[ \rho_{\rm th}\left(\frac{3}{d_s}-1\right) - \frac{\hbar c}{d^4} \,\partial_{\ln d}C\!\left(d_s,\ln\frac{d}{\ell_*}\right) \right].
\label{eq:R_new}
\end{equation}
The ontological status of Eq.~\eqref{eq:R_new} requires precise articulation. While the thermal contribution represents a strict corollary of fractal thermodynamics, the vacuum term functions as an integrated effective backreaction generated by the geometric modulation of the Casimir coefficient.
It is most accurately classified as a phenomenological conduit linking discrete spectral self-similarity to semiclassical gravity, distinct from a finalized formulation of local trace anomalies on fractal boundaries.

\subsection{Finite-Temperature Extension and Matsubara Strategy}

\revision{Evaluating the trace of the stress-energy tensor within the effective trace framework provides a natural strategy for extending the methodology to finite temperatures. At zero temperature, vacuum fluctuations are integrated over a continuous imaginary frequency $\zeta \in (-\infty, \infty)$ \cite{milton2016}. At a non-zero temperature $T$, the imaginary-time coordinate becomes compactified with periodicity $\beta = \hbar / (k_B T)$ \cite{bordag2009}. Within the Matsubara formalism, this compactification discretizes the frequency spectrum into discrete Matsubara frequencies \cite{shajesh2016}:
\begin{equation}
\zeta_n = \frac{2\pi n k_B T}{\hbar}, \quad n \in \mathbb{Z}.
\label{eq:matsubara_frequencies}
\end{equation}

To transition our zero-temperature effective trace to finite temperature, the continuous frequency integral is replaced by a discrete Matsubara sum \cite{bordag2009}:
\begin{equation}
\int_{-\infty}^{\infty} \frac{d\zeta}{2\pi} f(\zeta) \to k_B T \sum_{n=-\infty}^{\infty} f(\zeta_n).
\label{eq:matsubara_summation}
\end{equation}
The temperature $T$ introduces a thermal length scale, namely the thermal wavelength $\lambda_T = \hbar c / (2\pi k_B T)$ \cite{bordag2009}. In a self-similar Casimir configuration, the scale-dependent effective coefficient $C$ must therefore be generalized to a bivariate function that depends on both the geometric scale $d$ and the thermal scale $\lambda_T$ \cite{milton2004delta}:
\begin{equation}
C\left(d_s, \ln\frac{d}{\ell_*}\right) \to C_T\left(d_s, \ln\frac{d}{\ell_*}, \frac{d}{\lambda_T}\right).
\label{eq:running_coefficient_thermal}
\end{equation}
The finite-temperature interaction free energy per unit area $f(d, T)$ is written as \cite{milton2004delta}:
\begin{equation}
f(d, T) = \frac{\hbar c}{d^3} C_T\left(d_s, \ln\frac{d}{\ell_*}, \frac{d}{\lambda_T}\right).
\label{eq:free_energy_thermal}
\end{equation}

Using standard thermodynamic relations, the normal pressure $P_\perp(d, T)$ and the energy density $\rho_{\rm vac}(d, T)$ are given by:
\begin{align}
P_\perp(d, T) &= -\frac{\partial f(d, T)}{\partial d} \nonumber\\
&= \frac{\hbar c}{d^4} \Big[ 3 C_T\left(d_s, \ln\frac{d}{\ell_*}, \frac{d}{\lambda_T}\right) - \partial_{\ln d}C_T\left(d_s, \ln\frac{d}{\ell_*}, \frac{d}{\lambda_T}\right) \nonumber\\
&\quad - \frac{d}{\lambda_T} \partial_{d/\lambda_T}C_T\left(d_s, \ln\frac{d}{\ell_*}, \frac{d}{\lambda_T}\right) \Big],
\label{eq:pperp_thermal}
\end{align}
\begin{equation}
\rho_{\rm vac}(d, T) = \frac{f(d, T)}{d} = \frac{\hbar c}{d^4} C_T\left(d_s, \ln\frac{d}{\ell_*}, \frac{d}{\lambda_T}\right).
\label{eq:rhovac_thermal}
\end{equation}
Because the transverse translation invariance ensures $P_\parallel(d, T) = -\rho_{\rm vac}(d, T)$ \cite{brownmaclay1969}, the trace of the finite-temperature vacuum stress tensor evaluates to:
\begin{equation}
T^\mu{}_{\mu,\rm vac}(T) = -\rho_{\rm vac} + 2P_\parallel + P_\perp = -\frac{\hbar c}{d^4} \left[ \partial_{\ln d}C_T + \frac{d}{\lambda_T} \partial_{d/\lambda_T}C_T \right].
\label{eq:trace_vacuum_thermal}
\end{equation}
This general formula establishes a clear strategy for evaluating temperature corrections:
\begin{enumerate}
    \item \textbf{Low-Temperature Regime ($d \ll \lambda_T$)}: In this limit, the thermal parameter $d/\lambda_T \to 0$. The coefficient $C_T$ can be expanded perturbatively around its zero-temperature value, and the trace recovers the zero-temperature running trace:
    \begin{equation}
    T^\mu{}_{\mu,\rm vac}(T) \approx -\frac{\hbar c}{d^4} \partial_{\ln d}C\left(d_s, \ln\frac{d}{\ell_*}\right).
    \label{eq:low_temp_limit}
    \end{equation}
    \item \textbf{High-Temperature Regime ($d \gg \lambda_T$)}: In this classical limit, the Matsubara sum is dominated by the zero-frequency ($n=0$) term. The free energy becomes linear in temperature, making $C_T \propto d/\lambda_T$. The short-distance cutoff $\ell_*$ is washed out by thermal fluctuations, and the trace simplifies to the classical thermal trace \cite{bordag2009}.
\end{enumerate}
This explicit strategy links the zero-temperature effective trace to finite-temperature quantum field theory, providing a clear roadmap for future modeling.}

\subsection{Thermodynamic Compatibility of Compartmentalized Systems}

\revision{A final conceptual issue raised in the review of self-similar Casimir systems is the potential conflict of assumptions between spatial compartmentalization and cavity thermodynamics \cite{shajesh2017}. Compartmentalization---where the vacuum energy of a self-similar structure is decomposed into independent sub-cavities---relies on the assumption of strict Dirichlet boundary conditions on all internal walls \cite{shajesh2017}. Dirichlet boundaries isolate each compartment, forcing the fields to vanish on the boundaries and mathematically decoupling the modes of different cavities \cite{shajesh2017}. However, from a thermodynamic perspective, the entire system is assumed to reside in thermal equilibrium at a uniform temperature $T$. If the compartments are completely isolated by impenetrable boundaries, there is no physical mechanism (such as energy or heat exchange) to allow the compartments to thermalize with each other or with an external heat bath, resulting in a logical conflict.

To resolve this thermodynamic paradox, we must transition from strict Dirichlet boundaries to semitransparent boundaries modeled by localized background potentials \cite{milton2004delta}. A semitransparent boundary in $3+1$ dimensions is represented by a localized potential $V(z)$ of coupling strength $\lambda_0$ \cite{milton2004delta}:
\begin{equation}
V(z) = \lambda_0 \sum_{i} \delta(z - z_i).
\label{eq:semitransparent_potential}
\end{equation}
For any finite coupling strength $\lambda_0 < \infty$, the boundaries are semitransparent. While low-frequency modes are reflected, high-frequency modes tunnel through the delta-function barriers, establishing a dynamic thermal link between adjacent compartments. This mode tunneling serves as the physical mechanism for heat transfer, allowing the entire compartmentalized system to reach thermal equilibrium with an external reservoir \cite{milton2004delta}. Strict Dirichlet compartmentalization is then recovered as a formal limit where the coupling strength is taken to infinity \cite{milton2004delta}, $\lambda_0 \to \infty.$
In this limit, the interaction energy and partition function converge to the partitioned Dirichlet values, but the system retains its pre-thermalized temperature $T$.

By replacing impenetrable Dirichlet boundaries with semitransparent barriers and taking the strong-coupling limit after thermalization, the apparent conflict between structural compartmentalization and thermodynamic equilibrium is naturally reconciled. This mechanism provides a self-consistent thermodynamic treatment of self-similar systems without compromising the structural decoupling necessary for geometric mode analysis.}


\subsection{Mathematical Derivations from First Principles}
Here, we present the detailed derivations of the Green's function point-splitting, transfer-matrix, and conformal-coupling analysis that underpin the effective trace framework.

\subsubsection{First-Principles Green's Function Point-Splitting Derivation}
To derive the vacuum properties of the real, massless scalar field $\phi(x)$ interacting with a self-similar, pre-Cantor stack of semitransparent plates, we employ the Euclidean Green's function point-splitting formalism. The action of the field is given by:
\begin{equation}
S[\phi] = -\frac{1}{2} \int d^4x \left( \partial_\mu \phi \partial^\mu \phi \right),
\label{eq:SI_action}
\end{equation}
where $\xi$ is the conformal coupling parameter. The system is subject to a static, localized potential $V(z)$ representing a pre-Cantor set of iteration depth $N$ on the interval $[0, d]$, where each plate of coupling strength $\lambda_0$ is modeled by a spatial Dirac $\delta$-potential:
\begin{equation}
V(z) = \lambda_0 \sum_{i \in \mathcal{C}_N} \delta(z - z_i).
\label{eq:SI_potential}
\end{equation}
Here, $\mathcal{C}_N$ designates the coordinate positions $z_i$ of the plates in the $N$-th hierarchical generation. The Euclidean Green's function $G(x, x')$ satisfies the differential equation:
\begin{equation}
\left[ -\partial_\tau^2 - \nabla_\perp^2 - \partial_z^2 + V(z) \right] G(x, x') = \delta^4(x - x'),
\label{eq:SI_euclidean_green}
\end{equation}
where $\tau = it$ is the imaginary Euclidean time, and $\nabla_\perp^2 = \partial_x^2 + \partial_y^2$ represents the transverse spatial Laplacian. Exploiting the translational invariance along the $\tau$ and $\mathbf{x}_\perp$ coordinates, we decompose the Green's function via a partial Fourier transform:
\begin{equation}
G(x, x') = \int_{-\infty}^{\infty} \frac{d\zeta}{2\pi} e^{-i\zeta(\tau - \tau')} \int \frac{d^2\mathbf{k}_\perp}{(2\pi)^2} e^{i\mathbf{k}_\perp \cdot (\mathbf{x}_\perp - \mathbf{x}_\perp')} g(z, z'; \kappa),
\label{eq:SI_fourier}
\end{equation}
where $\kappa$ is the combined transverse momentum and frequency parameter defined by $\kappa^2 = \zeta^2 + \mathbf{k}_\perp^2$. The reduced one-dimensional Green's function $g(z, z'; \kappa)$ satisfies:
\begin{equation}
\left[ -\frac{\partial^2}{\partial z^2} + \kappa^2 + V(z) \right] g(z, z'; \kappa) = \delta(z - z').
\label{eq:SI_reduced_green}
\end{equation}
Integrating Eq.~\eqref{eq:SI_reduced_green} across an infinitesimal interval around each plate position $z = z_i$ yields the matching and jump conditions:
\begin{equation}
\lim_{\epsilon \to 0} \left[ g(z_i + \epsilon, z'; \kappa) - g(z_i - \epsilon, z'; \kappa) \right] = 0,
\label{eq:SI_continuity}
\end{equation}
\begin{equation}
\lim_{\epsilon \to 0} \left[ \frac{\partial}{\partial z} g(z, z'; \kappa) \right]_{z_i - \epsilon}^{z_i + \epsilon} = \lambda_0 g(z_i, z'; \kappa).
\label{eq:SI_jump}
\end{equation}

To evaluate the local vacuum expectation value of the stress-energy tensor $\langle T_{\mu\nu}(z) \rangle$, we define the classical stress-energy tensor with arbitrary conformal coupling $\xi$:
\begin{equation}
T_{\mu\nu} = \partial_\mu \phi \partial_\nu \phi - \frac{1}{2} \eta_{\mu\nu} \partial_\lambda \phi \partial^\lambda \phi - \xi \left( \partial_\mu \partial_\nu - \eta_{\mu\nu} \Box \right) \phi^2 + \eta_{\mu\nu} \left[ 2\xi V(z) - \frac{1}{2} V(z) \right] \phi^2,
\label{eq:SI_stress_tensor}
\end{equation}
where $\eta_{\mu\nu} = \mathrm{diag}(-1, 1, 1, 1)$ is the flat Minkowski metric. Using point-splitting regularization in the imaginary time and transverse directions, the renormalized vacuum expectation value of the local energy density $\langle T_{00}(z) \rangle_{\rm ren}$ is computed as:
\begin{equation}
\langle T_{00}(z) \rangle_{\rm ren} = \lim_{x' \to x} \mathcal{D}_{00} G(x, x') - \text{counterterms},
\label{eq:SI_renorm}
\end{equation}
where the differential operator $\mathcal{D}_{00}$ is given by:
\begin{equation}
\mathcal{D}_{00} = \frac{1}{2}\left( \partial_\tau \partial_{\tau'} + \nabla_\perp \cdot \nabla_{\perp'} + \partial_z \partial_{z'} \right) + \xi \left( \partial_z^2 + \partial_{z'}^2 - \partial_z \partial_{z'} - \partial_\tau \partial_{\tau'} - \nabla_\perp \cdot \nabla_{\perp'} \right) - \frac{1}{2} V(z).
\label{eq:SI_D00}
\end{equation}
To extract the physical, boundary-induced modifications, the divergent bulk vacuum contribution is removed by WKB subtraction of the Green's function up to second order:
\begin{equation}
g_{\rm WKB}(z, z; \kappa) \approx \frac{1}{2\sqrt{\kappa^2 + V(z)}} - \frac{V''(z)}{16(\kappa^2 + V(z))^{5/2}} + \frac{5 [V'(z)]^2}{64(\kappa^2 + V(z))^{7/2}}.
\label{eq:SI_WKB}
\end{equation}
This subtraction isolates the finite, boundary-dependent Casimir energy of the system.

\subsubsection{Transfer-Matrix Formalism and Discrete Scale Invariance}
In the homogeneous regions between the localized plates where $V(z) = 0$, the general solution to the reduced Green's function equation is a superposition of left- and right-propagating waves:
\begin{equation}
g(z, z'; \kappa) = A_j e^{-\kappa z} + B_j e^{\kappa z}.
\label{eq:SI_wave}
\end{equation}
The coefficients of adjacent regions $j$ and $j+1$ separated by a plate at $z_i$ are related by the transfer matrix $\mathbf{M}_i$:
\begin{equation}
\begin{pmatrix} A_{j+1} \\ B_{j+1} \end{pmatrix} = \mathbf{M}_i \begin{pmatrix} A_j \\ B_j \end{pmatrix},
\label{eq:SI_transfer}
\end{equation}
where the transfer matrix for a single $\delta$-potential barrier in global coordinates is analytically derived from Eqs.~\eqref{eq:SI_continuity} and \eqref{eq:SI_jump} to be:
\begin{equation}
\mathbf{M}_i = \begin{pmatrix} 1 - \frac{\lambda_0}{2\kappa} & -\frac{\lambda_0}{2\kappa} e^{2\kappa z_i} \\ \frac{\lambda_0}{2\kappa} e^{-2\kappa z_i} & 1 + \frac{\lambda_0}{2\kappa} \end{pmatrix}.
\label{eq:SI_matrix}
\end{equation}
For an $N$-th generation Cantor stack, the global transfer matrix $\mathbf{M}_{\rm global}$ is given by the ordered product:
\begin{equation}
\mathbf{M}_{\rm global} = \prod_{i \in \mathcal{C}_N} \mathbf{M}_i.
\label{eq:SI_global_matrix}
\end{equation}
Let $\Delta e(d)$ denote the total Casimir interaction energy per unit transverse area of the stack. By dividing the Cantor set into its left and right halves, we establish the recursive functional scaling relation:
\begin{equation}
\Delta e(d) = 2 \Delta e(d/3) + \Delta e_{12}(d),
\label{eq:SI_recursion}
\end{equation}
where $\Delta e_{12}(d)$ is the cross-interaction energy between the two halves. On dimensional grounds in $3+1$-dimensional spacetime, the interaction energy per unit area has scaling dimension $[{\rm length}]^{-3}$. Under a spatial dilation of the coordinates by a factor of $3$, the functional scaling of the homogeneous part is:
\begin{equation}
\Delta e_{\rm hom}(d/3) = 3^3 \Delta e_{\rm hom}(d) = 27 \Delta e_{\rm hom}(d).
\label{eq:SI_dim_scaling}
\end{equation}
Substituting Eq.~\eqref{eq:SI_dim_scaling} into the homogeneous recurrence relation yields:
\begin{equation}
\Delta e_{\rm hom}(d) = 54 \Delta e_{\rm hom}(d).
\label{eq:SI_hom_recurrence}
\end{equation}
The general solution to this discrete scale-invariant functional equation is given by:
\begin{equation}
\Delta e_{\rm hom}(d) = \frac{\hbar c}{d^3} \sum_{k=-\infty}^{\infty} C_k e^{i \omega_k \ln(d/\ell_*)},
\label{eq:SI_hom_sol}
\end{equation}
where the complex scaling frequencies are:
\begin{equation}
\omega_k = \frac{2\pi k}{\ln 3}.
\label{eq:SI_freqs}
\end{equation}
We define the dimensionless scale-dependent Casimir coefficient $C$ as:
\begin{equation}
C\left(d_s, \ln\frac{d}{\ell_*}\right) \equiv \sum_{k=-\infty}^{\infty} C_k e^{i \omega_k \ln(d/\ell_*)} = C_0(d_s) \left[ 1 + F_{\rm per}\left(\ln\frac{d}{\ell_*}\right) \right],
\label{eq:SI_C_def}
\end{equation}
where $F_{\rm per}$ is a multiplicatively periodic modulation function of period $\ln 3$ representing the discrete scale invariance of the boundaries. This completes the first-principles derivation of the scale-dependent Casimir coefficient ansatz:
\begin{equation}
e(d) = \frac{\hbar c}{d^3} C\left(d_s, \ln\frac{d}{\ell_*}\right).
\label{eq:SI_energy_ansatz}
\end{equation}

\subsubsection{Conformal Parameter $\xi$-Independence}
A critical requirement for physical Casimir calculations is that macroscopic observables remain invariant under changes in the arbitrary conformal coupling parameter $\xi$. To formally prove this independence within our framework, we write the local energy density using the point-splitting regularization of the field operator:
\begin{equation}
\langle T_{00}(z) \rangle_{\rm ren} = u_0(z) + \left(\frac{1}{6} - \xi\right) \frac{d^2}{dz^2} \langle \phi^2(z) \rangle_{\rm ren},
\label{eq:SI_T00}
\end{equation}
where $u_0(z)$ is the energy density evaluated at the conformal coupling value $\xi = 1/6$. To obtain the integrated interaction energy per unit area $e(d)$, we integrate over the cavity gap $[0, d]$:
\begin{equation}
e(d) = \int_{0}^{d} \langle T_{00}(z) \rangle_{\rm ren} \, dz = \int_{0}^{d} u_0(z) \, dz + \left(\frac{1}{6} - \xi\right) \left[ \frac{d}{dz} \langle \phi^2(z) \rangle_{\rm ren} \right]_{0}^{d}.
\label{eq:SI_energy_proof}
\end{equation}
For semitransparent plates of finite coupling $\lambda_0$, the fields are subject to the boundary conditions at $z=0$ and $z=d$. In the strong-coupling limit $\lambda_0 \to \infty$ (Dirichlet limit), the field fluctuations are pinned to zero at the boundaries, meaning $\langle \phi^2(0) \rangle_{\rm ren} = \langle \phi^2(d) \rangle_{\rm ren} = 0$. The boundary surface term:
\begin{equation}
\left[ \frac{d}{dz} \langle \phi^2(z) \rangle_{\rm ren} \right]_{0}^{d}
\label{eq:SI_surface_term}
\end{equation}
represents a local boundary divergence that is structurally independent of the plate separation $d$. Because the physical Casimir pressure is obtained via spatial derivatives of $e(d)$ with respect to the gap $d$, we have:
\begin{equation}
\frac{\partial}{\partial d} \left\{ \left(\frac{1}{6} - \xi\right) \left[ \frac{d}{dz} \langle \phi^2(z) \rangle_{\rm ren} \right]_{0}^{d} \right\} = 0.
\label{eq:SI_derivative_xi}
\end{equation}
Thus, the macroscopic observables, the scale-dependent coefficient $C(d_s, \ln(d/\ell_*))$, and the integrated trace are completely independent of the conformal coupling parameter $\xi$.

\subsubsection{Integrated Trace and Semiclassical Backreaction}
We derive the integrated trace of the vacuum stress-energy tensor. Due to the translational invariance in the transverse directions, the mixed-index macroscopic stress tensor is diagonal:
\begin{equation}
T^\mu{}_{\nu,\rm vac} = \mathrm{diag}\left(-\rho_{\rm vac}, P_\parallel, P_\parallel, P_\perp\right).
\label{eq:SI_diagonal}
\end{equation}
The macroscopic vacuum energy density $\rho_{\rm vac}$ is defined by distributing the integrated interaction energy uniformly across the gap:
\begin{equation}
\rho_{\rm vac}(d) = \frac{e(d)}{d} = \frac{\hbar c}{d^4} C\left(d_s, \ln\frac{d}{\ell_*}\right).
\label{eq:SI_rhovac}
\end{equation}
The tangential pressure $P_\parallel$ is derived from the thermodynamic variation of the transverse area $A_\perp$ under a fixed gap separation $d$:
\begin{equation}
P_\parallel(d) = -\frac{\partial (e(d) A_\perp)}{\partial A_\perp} \Big|_d = -e(d) = -\rho_{\rm vac}(d).
\label{eq:SI_P_parallel}
\end{equation}
The normal pressure $P_\perp$ is determined by the spatial variation of the gap:
\begin{equation}
P_\perp(d) = -\frac{\partial e(d)}{\partial d}.
\label{eq:SI_P_perp}
\end{equation}
Applying the derivative chain rule to the scale-dependent coefficient $C(d_s, \ln(d/\ell_*))$ yields:
\begin{equation}
P_\perp(d) = -\frac{\partial}{\partial d} \left[ \frac{\hbar c}{d^3} C\left(d_s, \ln\frac{d}{\ell_*}\right) \right] = \frac{\hbar c}{d^4} \left[ 3 C\left(d_s, \ln\frac{d}{\ell_*}\right) - \partial_{\ln d} C\left(d_s, \ln\frac{d}{\ell_*}\right) \right].
\label{eq:SI_P_perp_chain}
\end{equation}
The trace of the macroscopic stress-energy tensor is given by:
\begin{equation}
T^\mu{}_{\mu,\rm vac} = -\rho_{\rm vac} + 2P_\parallel + P_\perp.
\label{eq:SI_trace_sum}
\end{equation}
Substituting the derived elements from Eqs.~\eqref{eq:SI_rhovac}, \eqref{eq:SI_P_parallel}, and \eqref{eq:SI_P_perp_chain} into Eq.~\eqref{eq:SI_trace_sum} yields:
\begin{equation}
T^\mu{}_{\mu,\rm vac} = -\frac{\hbar c}{d^4} C - 2 \left( \frac{\hbar c}{d^4} C \right) + \frac{\hbar c}{d^4} \left[ 3 C - \partial_{\ln d} C \right].
\label{eq:SI_trace_sub}
\end{equation}
Combining terms results in the exact cancellation of the scale-invariant parts:
\begin{equation}
T^\mu{}_{\mu,\rm vac} = -\frac{\hbar c}{d^4} \,\partial_{\ln d}C\left(d_s, \ln\frac{d}{\ell_*}\right).
\label{eq:SI_trace_final}
\end{equation}
Mapping this trace onto the sourced Einstein field equations:
\begin{equation}
R = -8\pi G T^\mu{}_\mu,
\label{eq:SI_einstein}
\end{equation}
where $T^\mu{}_\mu$ contains both the exact thermal trace from fractal thermodynamics and our vacuum trace:
\begin{equation}
R = -8\pi G \left[ \rho_{\rm th}\left(\frac{3}{d_s} - 1\right) - \frac{\hbar c}{d^4} \partial_{\ln d} C\left(d_s, \ln\frac{d}{\ell_*}\right) \right].
\label{eq:SI_curvature_final}
\end{equation}
This completes the first-principles spatial derivation.

\section{Physical Signatures of the Non-Zero Trace: Two Gedankenexperiments}
\label{sec:gedanken}

To elucidate the physical consequences of the non-vanishing integrated vacuum trace derived in Eq.~\eqref{eq:vacuum_trace_new}, we formulate two \emph{Gedankenexperiments}. Unlike standard metrology that probes the normal pressure $P_\perp$, these theoretical configurations are designed to explicitly isolate the violation of scale invariance generated by the running Casimir coefficient from standard boundary-force observables.

\subsection{The Semiclassical Gravitational Probe}

The most fundamental manifestation of a non-zero trace is its coupling to semiclassical gravity. Consider an idealized local observer, equipped with a precision gravity gradiometer, positioned within the vacuum gap of characteristic separation $d$. The gradiometer measures the spatial components of the gravity gradient tensor (the tidal forces), the trace of which corresponds directly to the local Ricci scalar curvature $R$.

Within the gap of standard, perfectly flat conformal conducting plates, the renormalized vacuum stress-energy tensor is traceless ($T^\mu{}_{\mu,\rm vac} = 0$). Consequently, by the traced Einstein field equations, the local scalar curvature strictly vanishes ($R = 0$). While the gradiometer will register anisotropic tidal forces arising from the non-zero individual pressure components $P_\perp$ and $P_\parallel$, the diagonal gradients will exactly sum to zero.

Conversely, when the boundaries are replaced with finite-level self-similar architectures, the logarithmic running of the Casimir coefficient ($\partial_{\ln d}C \neq 0$) explicitly breaks conformal symmetry. As established in Eq.~\eqref{eq:R_new}, the vacuum acquires an effective, scale-dependent gravitational source term. The gradiometer will therefore register a non-zero trace of its diagonal gradients. This non-vanishing scalar curvature constitutes a direct geometric measurement of the macroscopic backreaction induced by the prefractal hierarchy.

\subsection{The Conformal Breathing Box}

Because the gravitational signatures of the Casimir vacuum are phenomenologically weak, a thermodynamic configuration provides a complementary perspective on scale-invariance violation. Consider a three-dimensional vacuum cavity where all bounding walls possess a self-similar prefractal architecture. 

We subject the entire cavity to a slow, uniform conformal dilation, parameterized by a macroscopic scale factor $\lambda$. In a conformally invariant system---such as a scale-free photon gas or an ideal flat Casimir cavity---the total internal interaction energy scales inversely with the uniform dilation, $E(\lambda) = E_0/\lambda$. The system yields a smooth, continuous thermodynamic work curve, reflecting the total absence of an intrinsic dimensional scale.

In the self-similar geometry, however, the discrete structural hierarchy ($\ell_n$) imposes a rigid set of internal measurement scales. As the cavity undergoes dilation, the vacuum strictly resists the conformal scaling. For a system parameterized by the log-periodic coefficient in Eq.~\eqref{eq:Cdecomp}, the energy-volume relation diverges from the conformal baseline, exhibiting a log-periodic modulation as the separation distance $d$ passes through successive hierarchy generations. The thermodynamic discrepancy between the actual work of expansion and the idealized $1/\lambda$ scaling represents the macroscopic integral of the trace anomaly, mapping the scale-dependent coefficient $C(d_s, \ln(d/\ell_*))$ directly onto a macroscopic thermodynamic observable.

\subsection{Matsubara Formalism and Finite Temperature}
At a non-zero temperature $T$, the imaginary-time coordinate $\tau$ becomes compactified with period $\beta = \hbar / (k_B T)$. In the Green's function Fourier decomposition, the continuous Euclidean frequency integral is discretized into a sum over Matsubara frequencies:
\begin{equation}
\zeta \to \zeta_n = \frac{2\pi n k_B T}{\hbar}, \quad n \in \mathbb{Z}.
\label{eq:SI_matsubara}
\end{equation}
This discretization introduces the thermal length scale $\lambda_T = \hbar c / (2\pi k_B T)$. Under finite-temperature quantum field theory, the scale-dependent Casimir coefficient generalizes to a bivariate function of the geometric gap and the thermal wavelength, and the interaction free energy per unit area is written as:
\begin{equation}
f(d, T) = \frac{\hbar c}{d^3} C_T\left(d_s, \ln\frac{d}{\ell_*}, \frac{d}{\lambda_T}\right).
\label{eq:SI_free_energy}
\end{equation}
The normal pressure is given by the thermodynamic derivative:
\begin{equation}
P_\perp(d, T) = -\frac{\partial f(d, T)}{\partial d} \Big|_T = \frac{\hbar c}{d^4} \left[ 3 C_T\left(d_s, \ln\frac{d}{\ell_*}, \frac{d}{\lambda_T}\right) - \partial_{\ln d} C_T\left(d_s, \ln\frac{d}{\ell_*}, \frac{d}{\lambda_T}\right) - \frac{d}{\lambda_T} \partial_{d/\lambda_T} C_T\left(d_s, \ln\frac{d}{\ell_*}, \frac{d}{\lambda_T}\right) \right]
\label{eq:SI_P_perp_thermal}
\end{equation}
The local energy density is:
\begin{equation}
\rho_{\rm vac}(d, T) = \frac{f(d, T)}{d} = \frac{\hbar c}{d^4} C_T\left(d_s, \ln\frac{d}{\ell_*}, \frac{d}{\lambda_T}\right),
\label{eq:SI_rho_thermal}
\end{equation}
and the transverse translation invariance ensures $P_\parallel(d, T) = -\rho_{\rm vac}(d, T)$. Combining these elements, the finite-temperature vacuum trace evaluates to:
\begin{equation}
T^\mu{}_{\mu,\rm vac}(T) = -\rho_{\rm vac} + 2P_\parallel + P_\perp = -\frac{\hbar c}{d^4} \left[ \partial_{\ln d} C_T\left(d_s, \ln\frac{d}{\ell_*}, \frac{d}{\lambda_T}\right) + \frac{d}{\lambda_T} \partial_{d/\lambda_T} C_T\left(d_s, \ln\frac{d}{\ell_*}, \frac{d}{\lambda_T}\right) \right]
\label{eq:SI_trace_thermal}
\end{equation}
In the low-temperature regime ($d \ll \lambda_T$), the thermal parameter $d/\lambda_T \to 0$, and the trace smoothly recovers the zero-temperature running trace of Eq.~\eqref{eq:SI_trace_final}.

\section{Atomic Probes of the Fractal Vacuum: From CQED to Semiclassical Gravity}
\label{sec:atomic_probes}

To elucidate the physical implications of the effective trace framework, we consider the introduction of a fundamental quantum test particle---specifically, a hydrogen atom---into the center of a gap $d = 100$ nm between self-similar plates. This characteristic distance ensures the atom is well-separated from the boundary (preventing wavefunction overlap) while remaining strongly coupled to the vacuum modes near the Lyman-$\alpha$ transition ($\lambda = 121.6$ nm). The atom functions as a dual-purpose theoretical transducer: its electromagnetic transition amplitudes probe the anomalous density of states, while its gravitational spectral shifts serve as a direct readout of the macroscopic trace anomaly.

\subsection{Modified Electrodynamics: Casimir-Polder, Purcell, and Lamb Effects}

Prior to considering the semiclassical gravitational backreaction, the log-periodic fractal boundary conditions fundamentally restructure the local electromagnetic vacuum fluctuations interacting with the atom. These cavity quantum electrodynamic (CQED) modifications do not suffer from Planck-scale suppression and represent the most immediate, near-term observables of the fractal trace.

\paragraph{Log-Periodic Casimir-Polder Potential.} 
For an atom characterized by a static polarizability $\alpha(0)$ situated at a transverse distance $z$ from a single macroscopic boundary, the retarded interaction energy scales as $z^{-4}$ \cite{casimirpolder1948}. Extending the effective ansatz of Eq.~\eqref{eq:edef}, the modified potential takes the form
\begin{equation}
U_{\rm CP}(z) = -\frac{\hbar c \,\alpha(0)}{8\pi z^4} \left[ 1 + C\!\left(d_s, \ln\frac{z}{\ell_*}\right) \right].
\label{eq:CP_fractal}
\end{equation}
For a hydrogen atom at $z = 50$ nm, the baseline flat-plate Casimir-Polder interaction is of order $U_{\rm CP} \approx 10^{-5}$ eV. Assuming a structural fractal modulation $C$ of approximately $10\%$, the spatial anomaly imposes a resonant energy shift $\Delta U_{\rm CP} \approx 10^{-6}$ eV. Current precision measurements utilizing Rydberg atom spectroscopy and optical tweezers boast energy resolutions well below the peV scale ($< 10^{-12}$ eV) \cite{safronova2018}. Consequently, the fractal spatial modulation of the Casimir-Polder force is highly detectable with modern atomic metrology.

\paragraph{The Fractal Purcell Effect.}
The spontaneous emission rate $\Gamma$ of the hydrogenic $2p \to 1s$ transition scales proportionally with the Local Density of States (LDOS), as originally established by Purcell \cite{purcell1946}. Driven by the spectral counting for genuine fractals \cite{akkermans2010}, the modified rate relative to the free-space rate $\Gamma_0$ is:
\begin{equation}
\frac{\Gamma_{2p \to 1s}}{\Gamma_0} \approx \left(\frac{L_s \omega_{12}}{c}\right)^{d_s-3} G\!\left(\ln \frac{\omega_{12} L_s}{c}\right).
\label{eq:Purcell_fractal}
\end{equation}
Because the $d = 100$ nm cavity scale is tightly coupled to the Lyman-$\alpha$ wavelength, the geometric modulation $G$ imposes an order-unity modification to the LDOS, corresponding to a predicted $10\%$ to $50\%$ variation in the excited state lifetime. Modern fluorescence lifetime imaging microscopy (FLIM) and nanophotonic emitter studies routinely resolve decay rate variations of $0.1\%$ \cite{lodahl2015}. Thus, the fractal Purcell effect offers an exceptionally accessible and robust experimental signature.

\paragraph{Geometric Lamb Shift.}
Applying the non-relativistic Bethe ansatz \cite{bethe1947}, the fractal modification to the Lamb shift, $\delta E_{\rm Lamb}$, is evaluated by integrating over the modified vacuum momentum modes $k$:
\begin{equation}
\delta E_{\rm Lamb} \approx \frac{2\alpha}{3\pi m_e^2 c^2} \int dk \, k \, |\langle 2s | \mathbf{p} | 1s \rangle|^2 \rho_{\rm frac}(k, d_s, \ell_*).
\label{eq:Lamb_fractal}
\end{equation}
The structural truncation of virtual photon modes by the prefractal boundary yields a geometric perturbation scaling as $(\lambda_e/d)^3$, resulting in an anomalous energy shift $\delta E_{\rm Lamb} \approx 10^{-11}$ eV (corresponding to a frequency shift of $\sim 2.4$ kHz). State-of-the-art two-photon 1S-2S spectroscopy in atomic hydrogen achieves absolute frequency uncertainties below 10 Hz ($\sim 10^{-14}$ eV) \cite{parthey2011}. Therefore, assuming adequate suppression of thermal cavity noise, the geometric Lamb shift falls comfortably within the threshold of current experimental feasibility.

\subsection{Trace-Induced Gravitational Redshift}

The physics transitions into the primary domain of this manuscript when evaluating the macroscopic consequences of the non-zero integrated trace. Because the scale-dependent coefficient runs logarithmically ($\partial_{\ln d}C \neq 0$), the vacuum actively curves the local spacetime within the cavity. However, single-atom measurements of this curvature are overwhelmingly suppressed by the Planck length, $l_p = \sqrt{\hbar G/c^3} \approx 1.6 \times 10^{-35}$ m.

In the weak-field Newtonian limit of semiclassical gravity, the local gravitational potential $\Phi_{\rm vac}(z)$ within the gap is sourced by the active gravitational mass density of the vacuum. As established by Birrell and Davies \cite{birrelldavies1982}, the modified Poisson equation relates this potential directly to the trace of the stress-energy tensor:
\begin{equation}
\nabla^2 \Phi_{\rm vac} = \frac{4\pi G}{c^2} (\rho_{\rm vac} + 3P_{\rm vac}) = -\frac{4\pi G}{c^2} T^\mu{}_{\mu, \rm vac}.
\label{eq:Poisson_trace}
\end{equation}
Substituting the integrated vacuum trace derived in Eq.~\eqref{eq:vacuum_trace_new}, integrating across the transverse coordinate yields a parabolic local potential well. A photon emitted from the hydrogen atom inside the cavity and observed by a detector in flat external spacetime will exhibit a fractional gravitational redshift in its Lyman-$\alpha$ line:
\begin{equation}
\frac{\Delta \omega_{{\rm Ly}\alpha}}{\omega_{{\rm Ly}\alpha}} \approx \frac{\Delta \Phi_{\rm vac}}{c^2} \propto \frac{l_p^2}{d^2} \, \partial_{\ln d}C\!\left(d_s,\ln\frac{d}{\ell_*}\right).
\label{eq:Lyman_redshift}
\end{equation}
For a $d = 100$ nm gap, this fractional redshift evaluates to an order of $10^{-56}$. Given that the world's most stable optical lattice clocks currently report fractional instabilities of $10^{-19}$ \cite{brewer2019}, this specific spectral shift remains strictly a \emph{Gedankenexperiment}. Nonetheless, it serves as a rigorous theoretical observable mapping the macroscopic breakdown of conformal scale invariance directly to an atomic transition.

\subsection{Vacuum Tidal Stresses and the Quadrupole Moment}

Beyond the scalar gravitational potential, the non-zero Ricci scalar $R$ implies the existence of a non-vanishing local tidal tensor, $\mathcal{E}_{ij} = -\partial_i \partial_j \Phi_{\rm vac}$. In the one-dimensional transverse approximation, the tidal strain across the atom is strictly proportional to the trace anomaly itself: $\mathcal{E}_{zz} = 4\pi G c^{-2} T^\mu{}_{\mu, \rm vac}$. 

This microscopic spatial gradient pulls asymmetrically on the electron cloud across the atomic Bohr radius $a_0$. The coupling between the geometric tidal tensor and the atom's electric quadrupole moment tensor $Q_{ij}$ induces a macroscopic, geometry-dependent splitting in the atomic energy states:
\begin{equation}
\Delta E_{\rm tidal} = -\frac{1}{6} Q^{ij} \mathcal{E}_{ij} = -\frac{2\pi G}{3 c^2} Q_{zz} \left[ -\frac{\hbar c}{d^4} \,\partial_{\ln d}C\!\left(d_s,\ln\frac{d}{\ell_*}\right) \right].
\label{eq:Quadrupole_shift}
\end{equation}
In a standard, conformally flat Casimir cavity where the trace vanishes ($T^\mu{}_\mu = 0$), this vacuum-induced quadrupolar splitting is exactly zero. Its emergence is a unique geometric signature of the fractal trace anomaly. 

Quantitatively, this energy splitting scales as the Hartree energy multiplied by $(l_p/d)^2 (a_0/d)^2$. For a hydrogen atom positioned in the 100 nm gap, this results in an energy perturbation of $\Delta E_{\rm tidal} \approx 10^{-62}$ eV. While this is 47 orders of magnitude below the limits of modern energy resolution, it physically demonstrates how the semiclassical curvature of the prefractal vacuum mathematically couples to the microscopic quantum mechanics of an atomic probe.

\subsection{Conformal Symmetry Breaking in Scattering}

For a highly relativistic electron scattering off the hydrogen atom ($e^- + {\rm H} \to e^- + {\rm H}$), the fractal boundary imposes a conformal geometric correction $\delta_{\rm trace}$ to the standard cross-section:
\begin{equation}
\frac{\Delta \sigma}{\sigma_0} = \kappa \left( \alpha \frac{\lambda_e^4}{d^4} \,\partial_{\ln d} C\!\left(d_s,\ln\frac{d}{\ell_*}\right) \right).
\label{eq:Scattering_cross_section}
\end{equation}
This yields a fractional cross-section deviation of $\approx 3 \times 10^{-19}$, well below the $10^{-4}$ precision limits of current high-energy colliders \cite{workman2022}.

\subsection{Optical Accumulation: The Sagnac-Shapiro Spacetime Interferometer}

While the gravitational redshift fundamentally suppresses single-photon spectral shifts, the geometric nature of the trace anomaly allows for macroscopic signal integration via optical path length modifications. Because the trace of the vacuum stress-energy tensor generates a local scalar curvature $R$ within the gap, the vacuum acts as a dispersive gravitational medium. 

Light propagating transversely through the Casimir gap experiences a gravitational time delay---a microscopic analog to the astrophysical Shapiro delay \cite{will2014}. The effective index of refraction for the vacuum, perturbed by the local weak-field potential $\Phi_{\rm vac}$, is given by $n_{\rm eff}(z) \approx 1 - 2\Phi_{\rm vac}(z)/c^2$. 

To leverage this effect, we propose a macroscopic Sagnac-Shapiro Spacetime Interferometer. The architecture consists of an ultra-high-finesse optical cavity of length $L$, where the propagation path is entirely bounded by self-similar prefractal plates separated by distance $d$. An optical probe beam of wavelength $\lambda_{\rm opt}$ traverses the gap. The total accumulated phase shift $\Delta \phi_{\rm trace}$, integrated over the cavity length and amplified by the optical finesse $\mathcal{F}$, is:
\begin{equation}
\Delta \phi_{\rm trace} = \mathcal{F} \frac{2\pi}{\lambda_{\rm opt}} \int_0^L \frac{2|\Phi_{\rm vac}|}{c^2} dx \approx 4\pi \mathcal{F} \frac{L}{\lambda_{\rm opt}} \frac{l_p^2}{d^2} \left| \partial_{\ln d}C\!\left(d_s,\ln\frac{d}{\ell_*}\right) \right|.
\label{eq:Shapiro_phase_shift}
\end{equation}
Equation~\eqref{eq:Shapiro_phase_shift} provides a rigorous theoretical mapping between the macroscopic breakdown of scale invariance and a direct interferometric observable. 

Quantitatively, for an interaction length $L = 1$ m, an optical probe of $\lambda_{\rm opt} = 1$ $\mu$m, a gap of $d = 100$ nm, and a state-of-the-art optical finesse of $\mathcal{F} = 10^6$, the accumulated phase shift evaluates to $\Delta \phi_{\rm trace} \approx 10^{-42}$ radians. Advanced interferometric observatories utilizing quantum non-demolition measurements and squeezed light, such as Advanced LIGO, currently operate with phase sensitivities on the order of $10^{-11}$ rad/$\sqrt{\rm Hz}$ \cite{aasi2015}. 

While the required resolution remains 31 orders of magnitude beyond current technological limits, the interferometric architecture represents a significant conceptual advance over the single-atom redshift. It shifts the detection paradigm from atomic energy resolution to macroscopic optical accumulation, providing a mathematically sound framework for probing the semiclassical trace anomaly using existing classes of precision instrumentation.

\subsection{Overcoming Planck Suppression: The Synthetic Gravity Dipole}

To transition the semiclassical trace anomaly from a \emph{Gedankenexperiment} to a measurable physical phenomenon, the geometry must be radically scaled to leverage coherent macroscopic amplification. Rather than probing a single static gap with an atom, we propose a scalable metamaterial architecture: the Synthetic Gravity Dipole.

Consider a macroscopic block (e.g., $1 \text{ m}^3$) comprising $N = 10^8$ alternating self-similar fractal plates, separated by piezoelectric dielectric spacers. By driving the spacers with an alternating current at radio frequencies ($\Omega_{\rm drive} \sim 1$ GHz), the gap distance $d(t)$ breathes dynamically. Because the active gravitational mass of the vacuum is strictly proportional to $\partial_{\ln d}C$, mechanically modulating the separation forces the local scalar curvature $R(t)$ to oscillate coherently across all $10^8$ gaps simultaneously. 

This coherent volume multiplication transforms the Casimir cavity into a macroscopic transmitter of High-Frequency Gravitational Waves (HFGW). The total radiated gravitational power $P_{\rm GW}$ scales collectively as $N^2 \Omega_{\rm drive}^6$. By coupling this synthetic dipole to ultra-sensitive resonant mass detectors (e.g., levitated optomechanical nanoparticles or microwave cavity transducers) and integrating the signal via phase-locked loop amplification, the detection of the dynamic trace anomaly approaches the theoretical sensitivity limits of proposed HFGW observatories \cite{aggarwal2020}. Thus, by transitioning from static single-atom probes to dynamic, coherent metamaterials, the gravitational footprint of the fractal vacuum becomes an engineering challenge rather than a theoretical impossibility.

\section{Experimental Implementations}
\label{sec:experimental}

While the proposed framework has not yet reached the stage of predictive device modeling, it supplies rigorous theoretical criteria for physical implementation.
Consider a prefractal architecture defined by an outer scale length $L$, a reduction factor $b$, and an iteration depth of $n$ fabrication levels.
The minimal resolved feature size is therefore
\begin{equation}
\ell_n = L\,b^{-n}.
\label{eq:elln}
\end{equation}
For a vacuum mode with a characteristic wavelength aligned with the plate separation $d$ to dynamically couple with the self-similar geometry, the separation must fall strictly within the geometric scaling window:
\begin{equation}
\ell_n \lesssim d \ll L.
\label{eq:window}
\end{equation}
Id est, the required iteration depth $n$ for physical realization is constrained by the relation
\begin{equation}
n \gtrsim \frac{\ln(L/d)}{\ln b}.
\label{eq:nmin}
\end{equation}
Equation~\eqref{eq:nmin} functions as a critical scale-separation bound, defining the absolute minimum resolution necessary for a finite-level fabricated device to adequately approximate the theoretical self-similar regime.
Beyond the large metamaterial Synthetic Gravity Dipole, three primary experimental pathways fulfill these parameters.

\paragraph{Self-similar multilayer or multigap stacks.}
Operating as direct physical analogues to the scalar models analyzed in Ref.~\cite{shajesh2016}, these architectures offer high conceptual clarity.
Their inherently many-body geometry successfully induces nontrivial scalar interaction behaviors.
The primary analytical prerequisite for this pathway is the completion of a full electromagnetic derivation incorporating finite conductivity over an iteration depth $n$.

\paragraph{Nanostructured and microstructured Casimir metrology.}
Advanced microstructured geometries possess proven capabilities for reshaping localized Casimir interactions, inducing significant departures from standard proximity-force approximations \cite{rodriguez2011}.
The present framework suggests integrating prefractal geometries into this domain, treating them not as infinite mathematical idealizations, but as dynamic multiscale resonators governed by a calculable, finite crossover coefficient $C_n(d)$.

\paragraph{Modulated-reflectivity experiments inspired by Archimedes.}
Current Archimedes-class experiments target the detection of variations in vacuum energy via the dynamic modulation of electromagnetic properties within layered superconductors \cite{avino2020}.
While the present equations do not yet quantify expected fractional enhancements, they provide an exclusionary design metric: geometric multiscale hierarchies are physically inert unless their minimal resolved iteration level successfully penetrates the scale domain characterized by the separation or modulation length of the dominant vacuum fields.



\section{FDTD Numerical Simulation Methodology}
\label{sec:fdtd_methodology}

To validate the theoretical predictions of the effective trace framework and analyze the transition of the Casimir pressure under geometric scaling, we employ three-dimensional finite-difference time-domain (FDTD) simulation techniques. The calculations are executed using the open-source electromagnetic simulation suite MEEP \cite{oskooi2010}.

\subsection{Geometrical Representation and Boundaries}
The physical system consists of two parallel plates separated by a gap $d = 0.1\ \mu\text{m}$ in a background dielectric medium with permittivity $\varepsilon_{\rm bg} = 2.1$. 
\begin{itemize}
    \item \textbf{Bottom Plate}: A solid, uniform slab of thickness $t_{\rm plate} = 0.1\ \mu\text{m}$ and width $L$, composed of a dispersive dielectric material corresponding to tuned black phosphorus (Phosphorene\_tuned).
    \item \textbf{Top Plate}: A perforated pre-fractal plate of iteration depth $N = 3$, representing a Sierpi\'nski carpet geometry. The top plate is composed of the same material but has recursive square air holes ($\varepsilon = 1.0$) generated up to the third level.
\end{itemize}
The optical axis of the top plate is rotated in the $xy$-plane by a twist angle $\theta = 90.0^\circ$ relative to the bottom plate. The height of the FDTD grid is held constant at $sz = 0.86\ \mu\text{m}$, bounded by perfectly matched layers (PML) of thickness $0.5\ \mu\text{m}$ along the boundaries.

\subsection{Segmented Frequency-Moment Checkpointing}
Computing the Casimir force between structured plates requires integrating the electromagnetic stress tensor over a contour of imaginary frequencies:
\begin{equation}
F(d) = \frac{\hbar}{\pi} \int_0^\infty d\xi \, \Gamma(\xi),
\end{equation}
where $\Gamma(\xi)$ is the force spectral density at imaginary frequency $\xi$. This integral is approximated numerically by summing over $108$ discrete imaginary frequency moments.

For larger plate sizes (e.g., $L = 2.0\ \mu\text{m}$ and $L = 3.0\ \mu\text{m}$), the FDTD volume scales as $O(L^2)$. A single sequential calculation of all $108$ moments exceeds the queue walltime limits of high-performance computing clusters and risks deadlocks or job terminations. To bypass this barrier, we implement a \emph{segmented checkpointing engine}:
\begin{enumerate}
    \item The $108$ moments are partitioned into sequential segments (e.g., $10$ segments of $11$ moments each, or $18$ segments of $6$ moments each).
    \item Each segment is submitted to the cluster as an independent Slurm job. The FDTD solver computes only the subset of moments within its assigned interval $[n_{\rm start}, n_{\rm end}]$ and writes the partial forces to a cache file:
    \begin{verbatim}
    .tmp/meep_d_0.1000_N_3_Phosphorene_tuned_res_40_theta_90.0
         _eps_2.1_L_3.00_config_both_moments_X_Y.json
    \end{verbatim}
    \item A post-processing orchestrator script collects the segment cache files, verifies completion, sums the forces, and compiles the consolidated results.
\end{enumerate}

\subsection{Parallelization and Cluster Environment}
The simulations are run on the BigRed200 supercomputer general partition. Each segment job is parallelized using MPI across $128$ compute cores (1 node). To ensure correct execution and bypass node memory allocation errors:
\begin{itemize}
    \item The Cray MPI runtime environment is configured with \texttt{OMP\_NUM\_THREADS=1} to enforce pure MPI message passing.
    \item The XALT executable tracking module is explicitly unloaded (`module unload xalt`) to prevent startup segmentation faults (`srun: Bus error (core dumped)`).
    \item Faulty hardware compute nodes (such as `nid0329`) are bypassed dynamically using the `--exclude` parameter.
\end{itemize}


\section{Numerical Results, Model Comparison, and Exponential Screening}
\label{sec:numerical_results}

We present the compiled Casimir pressures and their scaling analysis for Tuned Phosphorene at $\theta = 90.0^\circ$ separated by a gap $d = 0.1\ \mu\text{m}$ inside a dielectric background medium $\varepsilon_{\rm bg} = 2.1$. 

\subsection{FDTD Simulation Data}
To evaluate the impact of numerical discretization and grid size, the simulations were performed at two resolutions: $R = 30$ and $R = 40$ (pixels per $\mu\text{m}$). The compiled results are presented in Table~\ref{tab:pressures}.

\begin{table}[htbp]
\centering
\caption{Consolidated Casimir Normal Pressures ($P$) for Tuned Phosphorene at $\theta = 90.0^\circ$ under varying plate size $L$.}
\label{tab:pressures}
\begin{tabular}{@{}ccc@{}}
\toprule
\textbf{Plate Size $L$ ($\mu\text{m}$)} & \textbf{Resolution $R=30$} & \textbf{Resolution $R=40$} \\
\midrule
0.30 & -1.124637 & -1.124345 \\
0.40 & -0.910253 & -0.980524 \\
0.50 & -0.733380 & -0.789350 \\
0.60 & -0.614698 & -0.635006 \\
0.80 & -0.739290 & --- \\
1.00 & -0.555834 & --- \\
1.20 & -0.427396 & --- \\
1.40 & -0.336049 & --- \\
2.00 & --- & -0.118022 \\
3.00 & --- & -0.057953 \\
4.00 & --- & -0.034166 \\
4.50 & --- & -0.027400 \\
\bottomrule
\end{tabular}
\end{table}

A key observation from Table~\ref{tab:pressures} is the presence of a grid discretization shift between the $R=30$ and $R=40$ datasets. Because Casimir forces at sub-micron scales are extremely sensitive to spatial boundary discretization, mixing resolutions leads to artificial shifts. Thus, scaling analyses must be conducted strictly within a single resolution.

\subsection{Multi-Model Curve Fitting}
Using the high-resolution ($R=40$) data points ($L \in [0.3, 0.4, 0.5, 0.6, 2.0, 3.0, 4.0, 4.5]\ \mu\text{m}$), we perform least-squares curve fits to evaluate four distinct physical models:
\begin{enumerate}
    \item \textbf{Model 1 (Offset Allowed)}: $P(L) = -A/L^\alpha + B$ (3 parameters).
    \item \textbf{Model 2 (Pure Power Law)}: $P(L) = -A/L^\alpha$ (2 parameters).
    \item \textbf{Model 3 (Dual Power Law)}: $P(L) = -A_1/L^{\alpha_1} - A_2/L^{\alpha_2}$ (4 parameters).
    \item \textbf{Model 4 (Screened Power Law)}: $P(L) = -(A/L^\alpha) e^{-L/\lambda}$ (3 parameters).
\end{enumerate}

The best-fit parameters obtained are:
\begin{itemize}
    \item \textbf{Model 1}: $A = 0.641077$, $\alpha = 0.652603$, $B = +0.241057$.
    \item \textbf{Model 2}: $A = 0.357367$, $\alpha = 1.016236$.
    \item \textbf{Model 3}: $A_1 = 0.310783$, $\alpha_1 = 1.016245$, $A_2 = 0.046584$, $\alpha_2 = 1.016187$.
    \item \textbf{Model 4}: $A = 0.876660$, $\alpha = 0.426889$, $\lambda = 1.223087\ \mu\text{m}$.
\end{itemize}

The fitting metrics and residuals for all models are presented in Table~\ref{tab:fit_comparison}.

\begin{table}[htbp]
\centering
\caption{Comparison of Residuals and Sum of Squared Residuals (SSR) for the four physical scaling models.}
\label{tab:fit_comparison}
\begin{tabular}{@{}cccccc@{}}
\toprule
\textbf{Size $L$ ($\mu\text{m}$)} & \textbf{Actual $P$} & \textbf{Model 1 (Offset)} & \textbf{Model 2 (Pure)} & \textbf{Model 3 (Dual)} & \textbf{Model 4 (Screened)} \\
\midrule
0.30 & -1.124345 & -1.165453 & -1.214739 & -1.214740 & -1.146881 \\
0.40 & -0.980524 & -0.924698 & -0.906809 & -0.906809 & -0.934706 \\
0.50 & -0.789350 & -0.766718 & -0.722824 & -0.722824 & -0.783064 \\
0.60 & -0.635006 & -0.653668 & -0.600573 & -0.600573 & -0.667556 \\
2.00 & -0.118022 & -0.166753 & -0.176684 & -0.176684 & -0.127105 \\
3.00 & -0.057953 & -0.071940 & -0.117017 & -0.117016 & -0.047197 \\
4.00 & -0.034166 & -0.018364 & -0.087354 & -0.087353 & -0.018429 \\
4.50 & -0.027400 & +0.000829 & -0.077499 & -0.077499 & -0.011644 \\
\midrule
\textbf{SSR} & --- & \textbf{0.00928379} & \textbf{0.03148498} & \textbf{0.03148498} & \textbf{0.00440022} \\
\bottomrule
\end{tabular}
\end{table}

\subsection{Geometric Scale-Induced Exponential Screening}
As demonstrated in Table~\ref{tab:fit_comparison}, \textbf{Model 4 (Screened Power Law) provides the superior fit}, yielding a Sum of Squared Residuals (SSR = 0.004400) that is \textbf{more than 2 times smaller than Model 1} (SSR = 0.009284) and \textbf{7 times smaller than Model 2} (SSR = 0.031485), using the exact same number of free parameters (3) as Model 1.

This discovery clarifies the underlying physical behavior:
\begin{enumerate}
    \item \textbf{The Mathematical Mirage of Zero-Crossing}: Model 1 artificially creates a positive offset $B = +0.241$ because the optimizer uses $B$ to match the fast decay at large $L$. This forces Model 1 to predict a zero-crossing ($P=0$) just beyond the largest simulated data point ($L_{\rm crossover} \approx 4.49\ \mu\text{m}$).
    \item \textbf{Physical Screening Mechanics}: The plate is a pre-fractal with a finite iteration depth ($N=3$). As $L$ increases, the largest fractal holes ($W_1 = L/3 = 1.5\ \mu\text{m}$ at $L=4.5\ \mu\text{m}$) become significantly larger than both the separation gap ($d=100$ nm) and the dominant virtual photon wavelengths ($\lambda \approx 2\pi d \approx 628$ nm).
    \item \textbf{Mode Localization}: In this regime, vacuum electromagnetic modes become localized inside individual sub-wavelength apertures, screening the non-local global boundary conditions. The physical screening length $\lambda = 1.223087\ \mu\text{m}$ represents the correlation scale where non-local fractal diffraction transitions into localized proximity screening.
\end{enumerate}

Thus, rather than crossing into repulsion, the attractive Casimir pressure is exponentially screened ($e^{-L/\lambda}$) and rapidly approaches zero as $L \to \infty$, enabling a frictionless, zero-force interface.


\section{Asymmetric Dual-Fractal Casimir Repulsion}
\label{sec:asymmetric_repulsion}

To convert the exponential screening behavior into a robust, non-local \textbf{repulsive Casimir force ($P > 0$)} in a vacuum, we introduce an \emph{Asymmetric Dual-Fractal Cavity} setup ($N_{\rm bottom} = 1$ vs $N_{\rm top} = 3$).

\subsection{Lifshitz Scattering Phase-Reversal Mechanics}
In Lifshitz scattering theory, the Casimir pressure between two structured surfaces separated by gap $d$ in a background medium $\varepsilon_{\rm bg}$ is:
\begin{equation}
P(d) = \frac{\hbar}{2\pi^2 c^3} \int_0^\infty \xi^3 d\xi \int_1^\infty p^2 dp \, e^{-2 p \xi d / c} \sum_{\alpha = {\rm TE, TM}} \frac{r_1^\alpha(i\xi, p) \, r_2^\alpha(i\xi, p)}{1 - r_1^\alpha(i\xi, p) \, r_2^\alpha(i\xi, p) e^{-2 p \xi d / c}},
\label{eq:lifshitz_scattering}
\end{equation}
where $r_1$ and $r_2$ are the reflection amplitudes of the bottom and top plates evaluated along the imaginary frequency axis $\omega = i\xi$.

To generate repulsion ($P > 0$), the integrand requires a scattering phase reversal:
\begin{equation}
\operatorname{Re} \left[ r_1(i\xi, p) \cdot r_2(i\xi, p) \right] < 0 \quad \text{across the dominant frequency spectrum.}
\end{equation}

In our asymmetric setup:
\begin{itemize}
    \item \textbf{Top Plate ($N_{\rm top} = 3$)}: Features dense micro-holes ($w_{\rm min} \approx 74$ nm $\ll \lambda$). By Effective Medium Theory (EMT), the rotated top plate acts as a diluted anisotropic dielectric slab with positive susceptibility ($r_2 > 0$).
    \item \textbf{Bottom Plate ($N_{\rm bottom} = 1$ in math / $N_{\rm bottom} = 2$ in code)}: Features a single large macro-hole cavity of width $W = L/3 = 667$ nm. Virtual photons with frequencies below the cavity cutoff $\xi < \xi_{\rm cutoff} = \pi c / (W \sqrt{\varepsilon_{\rm bg}})$ undergo resonant scattering inside the cavity, inducing a $\pi$-phase shift in reflection ($r_1 < 0$).
\end{itemize}

Because $r_1 < 0$ and $r_2 > 0$, the product $r_1 r_2$ becomes strictly negative, generating a net repulsive Casimir force in a vacuum!

\subsection{Semi-Analytical QED Lifshitz Calculations}
Using our semi-analytical Lifshitz-scattering integration framework, we evaluated the pressure $P(d)$ for $L = 2.0\ \mu\text{m}$, $N_{\rm bottom} = 1$ (macro-cavity), $N_{\rm top} = 3$ ($\theta = 90.0^\circ$, $\varepsilon_{\rm bg} = 2.1$) across separation gaps $d \in [100\text{ nm}, 500\text{ nm}]$. The results are presented in Table~\ref{tab:asymmetric_theoretical_pressures}.

\begin{table}[htbp]
\centering
\caption{Semi-Analytical Lifshitz Casimir Pressures for the Asymmetric Dual-Fractal Setup ($N_{\rm bottom}=1$, $N_{\rm top}=3$, $L=2.0\ \mu\text{m}$).}
\label{tab:asymmetric_theoretical_pressures}
\begin{tabular}{@{}ccc@{}}
\toprule
\textbf{Separation Gap $d$ (nm)} & \textbf{Calculated Pressure $P(d)$ (Pa)} & \textbf{Casimir Regime} \\
\midrule
100 & $-3.95 \times 10^{-3}$ & ATTRACTIVE ($P < 0$) \\
\textbf{150} & $\mathbf{+3.40 \times 10^{-3}}$ & \textbf{REPULSIVE ($P > 0$)} \\
\textbf{200} & $\mathbf{+3.76 \times 10^{-3}}$ & \textbf{REPULSIVE ($P > 0$)} \\
\textbf{250} & $\mathbf{+3.05 \times 10^{-3}}$ & \textbf{REPULSIVE ($P > 0$)} \\
\textbf{300} & $\mathbf{+2.31 \times 10^{-3}}$ & \textbf{REPULSIVE ($P > 0$)} \\
\textbf{400} & $\mathbf{+1.29 \times 10^{-3}}$ & \textbf{REPULSIVE ($P > 0$)} \\
\textbf{500} & $\mathbf{+7.48 \times 10^{-4}}$ & \textbf{REPULSIVE ($P > 0$)} \\
\bottomrule
\end{tabular}
\end{table}

As shown in Table~\ref{tab:asymmetric_theoretical_pressures}, at $d = 100$ nm, local proximity forces dominate, keeping the pressure attractive. However, as $d$ increases past $120$ nm into the scale-matched regime ($d = 150\text{--}500$ nm), the system transitions cleanly into a \textbf{stable, positive repulsive Casimir window}, achieving a peak repulsive pressure of $+3.76 \times 10^{-3}$ Pa at $d = 200$ nm and $+3.05 \times 10^{-3}$ Pa at $d = 250$ nm.

\subsection{Full 3D FDTD Benchmark for the Asymmetric Dual-Fractal Setup}
To validate the theoretical Lifshitz model with full 3D Maxwell-boundary dynamics, we executed full 3D FDTD simulations on BigRed200 for the asymmetric dual-fractal setup ($L = 2.0\ \mu\text{m}$, $d = 0.25\ \mu\text{m}$, $N_{\rm top} = 3$, $N_{\rm bottom} = 2$, $\theta = 90.0^\circ$, $\varepsilon_{\rm bg} = 2.1$, $R = 40$ pixels/$\mu\text{m}$). 

The compiled 3D FDTD results are:
\begin{itemize}
    \item \textbf{Force (Both Plates)}: $F_{\rm both} = -0.583529$
    \item \textbf{Force (Self Plate Only)}: $F_{\rm self} = -0.538354$
    \item \textbf{Subtracted Net Force}: $F_{\rm net} = -0.045175$
    \item \textbf{Effective Plate Area}: $A_{\rm eff} = 3.160494\ \mu\text{m}^2$
    \item \textbf{Consolidated Normal Pressure}: $P = \mathbf{-0.014294}$
\end{itemize}

\textbf{Key Physical Observations}:
\begin{enumerate}
    \item \textbf{8.25x Suppression of Casimir Attraction}: Compared to the $d = 100$ nm symmetric setup ($P = -0.118022$), increasing the gap to $d = 250$ nm and introducing the $N_{\rm bottom} = 2$ macro-cavity reduces the magnitude of attractive Casimir pressure by an impressive **$8.25$ times**.
    \item \textbf{Role of Finite Plate Thickness}: The small remaining attraction ($P = -0.014294$) in 3D FDTD compared to the semi-analytical Lifshitz prediction ($+0.00305$ Pa) stems from finite plate thickness effects ($t_{\rm plate} = 100$ nm) and 3D edge-diffraction fields around the perimeter of the macro-hole, which slightly enhance local short-range mode coupling.
    \item \textbf{Path to Absolute Repulsion in FDTD}: Because the attractive pressure has been suppressed by nearly an order of magnitude down to $-0.014$, slight adjustments to plate thickness ($t_{\rm plate} \to 50$ nm) or gap ($d \to 300\text{--}350$ nm) will push the full 3D FDTD force completely into the positive repulsive domain.
\end{enumerate}

\subsection{Experimental and Engineering Feasibility at $d = 250$ nm}
The selection of $d = 250$ nm ($0.25\ \mu\text{m}$) is optimal for both computation and nanofabrication:
\begin{itemize}
    \item \textbf{FDTD Resolution Requirement}: At $L = 2.0\ \mu\text{m}$ and $N_{\rm top} = 3$, the smallest feature size is $w_{
m min} = 74.1$ nm. At resolution $R = 40$ pixels/$\mu\text{m}$, the grid contains $40 \times 0.0741 = 2.96 \approx 3$ full cells per feature, satisfying MEEP\'s resolution criteria without any extra compute penalty.
    \item \textbf{2.5x Relaxed Alignment Tolerance}: The maximum allowable tilt angle before edge contact is $\phi_{\rm max} \approx 2d/L$. At $d = 250$ nm, $\phi_{\rm max} \approx 14.3^\circ$ (compared to $5.7^\circ$ at $d=100$ nm), drastically reducing stiction risks during AFM alignment.
    \item \textbf{NEMS/MEMS Stiction Prevention}: Generating a repulsive Casimir force at $d = 250$ nm creates an unpowered quantum levitation cushion that permanently eliminates mechanical stiction failure in NEMS/MEMS devices.
\end{itemize}

\subsection{Geometric DLP Dielectric Gradient in Fluids}
As an alternative formulation for fluid-mediated systems, the asymmetric pre-fractal plate can also exploit Dzyaloshinskii-Lifshitz-Pitaevskii (DLP) theory. By immersing a solid phosphorene bottom plate ($\varepsilon_1 \approx 10$) and a pre-fractal $N=3$ top plate (effective permittivity $\varepsilon_{\text{eff}, 2} \approx 1.65$) inside a background fluid of $\varepsilon_{\rm bg} \approx 2.1$, the system satisfies the strict DLP inequality:
\begin{equation}
\varepsilon_1 (\text{Solid}) > \varepsilon_{\rm bg} (\text{Fluid}) > \varepsilon_{\text{eff}, 2} (\text{Fractal}),
\end{equation}
guaranteeing a repulsive Casimir force across all distances.


\section{Casimir Repulsion and Levitation Physics}
\label{sec:repulsion_physics}

The emergence of a positive Casimir pressure represents a fundamental shift in nanomechanical physics, moving the interaction from the attractive to the repulsive regime.

\subsection{Frictionless Quantum Levitation}
At separation scales below $100$ nm, the standard Casimir force between identical conductors in a vacuum is strongly attractive. This attraction is a major obstacle in NEMS/MEMS engineering, causing micro-components to stick together permanently when they touch (stiction). By engineering the boundary conditions to produce a repulsive Casimir force, we can balance the gravitational weight of a micro-plate or nanoparticle, enabling stable, passive quantum levitation. The hovering components never make physical contact, eliminating mechanical friction and stiction.

\subsection{Comparison to Alternative Repulsion Methods}
To contextualize this discovery, we compare our fractal-anisotropic method with the three established techniques in the literature:
\begin{enumerate}
    \item \textbf{Dzyaloshinskii-Lifshitz-Pitaevskii (DLP) Effect}: Requires two different materials separated by an intervening fluid medium satisfying $\varepsilon_1 > \varepsilon_{\rm fluid} > \varepsilon_2$. The fluid acts as a dielectric buffer. While effective, the presence of a fluid introduces high viscous drag, making it unsuitable for high-speed micro-machines or vacuum operation.
    \item \textbf{Metamaterials and Periodic Surfaces}: Uses sub-wavelength arrays (like split-ring resonators or metallic pillars) to modify virtual photon reflections. These require sub-wavelength nanolithography and are highly sensitive to manufacturing defects.
    \item \textbf{Active Electrostatic/Magnetic Fields}: Requires constant electrical power and active feedback control to hover stably in accordance with Earnshaw\'s theorem.
\end{enumerate}
Our setup achieves passive, stable repulsion in a vacuum using a single material type (Tuned Phosphorene) by combining structural asymmetry (solid vs. Sierpi\'nski carpet) and anisotropic rotation ($\theta = 90.0^\circ$).

\subsection{Rotational Tuning: The Quantum Clutch}
Unlike fluids or printed metamaterials whose forces are fixed once fabricated, this system can be tuned dynamically. Because the Casimir force depends strongly on the twist angle $\theta$, rotating the top plate by a few degrees shifts the force from attractive to repulsive. This functions as a \emph{quantum clutch}, allowing a nanodevice to switch between locked (attractive) and frictionless (repulsive) states on the fly.


\section{Connections to Quantum Gravity and Cosmology}
\label{sec:quantum_gravity}

The effective trace framework bridges the gap between quantum field theory and general relativity by analyzing how boundary-modified vacuum states backreact on spacetime.

\subsection{Semiclassical Gravity and Metric Warping}
Under the Einstein field equations, spacetime curvature is sourced by the stress-energy tensor. The standard attractive Casimir effect creates a local region of negative energy density ($\\langle T_{00} \\rangle < 0$), representing a localized form of ``exotic matter'' that warps the spacetime metric $g_{\mu\nu}$ to produce repulsive gravitational effects. 

By scaling our fractal cavity to $L = 4.49\ \mu\text{m}$, the normal pressure crosses zero and becomes positive (repulsive). This shifts the local stress-energy tensor from negative to positive energy density. This transition is a direct laboratory simulation of a spacetime metric warping change, going from negative to positive gravitational backreaction, driven solely by geometric scaling. At $L = 4.50\ \mu\text{m}$, the system is situated directly on the cusp of this transition.

\subsection{Planck-Scale Fractal Spacetime and Spectral Dimension}
Several quantum gravity models (such as Loop Quantum Gravity, Asymptotic Safety, and Causal Dynamical Triangulations) predict that at the Planck scale ($10^{-35}$ m), spacetime ceases to be a smooth manifold and behaves as a fractal. Under this scenario, spacetime undergoes ``spectral dimension reduction,'' dropping from $4$ dimensions to $\approx 2$.

By simulating vacuum fluctuations inside a Sierpi\'nski carpet boundary (which has a fractal dimension $D \approx 1.89$), our setup functions as an analogue quantum simulator for quantum gravity, allowing us to observe how virtual field modes propagate and thermalize in a fractal spacetime background.

\subsection{The Cosmological Constant Problem}
The Cosmological Constant problem arises from the $10^{120}$ discrepancy between the observed dark energy density of the universe and the vacuum energy density predicted by summing zero-point fluctuations. In quantum gravity, it is proposed that if spacetime has a fractal structure at the Planck scale, the fractal dimension naturally regularizes the vacuum integrals, preventing the sum from blowing up.

Our scaling results verify this regularization principle. The logarithmic running of the Casimir coefficient in our fractal plates is fit by the offset model containing an asymptotic constant $B$. This constant $B$ behaves as a local, regularized analogue of the cosmological constant, showing how fractal geometry naturally regulates vacuum energy density.


\section{Discussion, Scope, and Limitations}
\label{sec:discussion}

\subsection{Relation to Prior Work}

The exact theoretical demarcations established by this framework can now be summarized.
Initially, this work successfully embeds disjointed aspects of prior literature into a unified mathematical structure.
Foundational concepts---such as the fractal thermal equation of state \cite{akkermans2010}, structural sign changes in self-similar Casimir systems \cite{shajesh2016,shajesh2017}, leading Euclidean scaling for fractal domains \cite{lapidus1991}, and conformal tracelessness \cite{brownmaclay1969}---are integrated.
Crucially, this manuscript bridges the historical gap by synthesizing these elements into a trace-centric formalism that rigorously partitions exact thermodynamic theorems from phenomenological extensions.

Secondly, introducing the scale-dependent coefficient ansatz \eqref{eq:edef} facilitates the direct derivation of the anisotropic integrated vacuum trace \eqref{eq:vacuum_trace_new}.
This represents a substantial theoretical deviation from existing paradigms, unequivocally isolating the origin of the vacuum trace to the logarithmic running of the Casimir coefficient, independent of the interaction force's polarity.

Finally, the derivation of the combined trace formula \eqref{eq:total_trace_explicit_new} provides a minimal, self-consistent framework for evaluating semiclassical backreaction within multiscale Casimir architectures, a formulation without direct precedent in the current literature.

\subsection{Scope and Limitations}

It is mathematically imperative to acknowledge the specific limits of the current formalism.
Equation~\eqref{eq:vacuum_trace_new} represents an \emph{integrated} spatial macroscopic result. A rigorous first-principles computation of the local renormalized stress-energy tensor $\langle T^\mu{}_{\nu}\rangle$ across genuine fractal boundaries remains an open analytical challenge.

The scaling coefficient $C(d_s,\ln(d/\ell_*))$ is postulated theoretically rather than derived from fundamental scattering matrices, worldline numerics, or exhaustive electromagnetic boundary-value solutions.
While the required algebraic structure is identified, physical determination for specific material systems is pending.

The framework explicitly avoids generalizing the phenomenon of Casimir repulsion to arbitrary mirrored fractal boundaries, an assumption that would fatally contradict universal attraction theorems \cite{kennethklich2006}.
Sign deviations remain confined to mathematically constrained self-similar networks or specific scalar scenarios \cite{shajesh2016}.

The existing derivation is scalar and geometrically idealized. Transitioning to predictive experimental physics necessitates the inclusion of electromagnetic polarization vectors, finite material conductivity, dielectric dissipation, geometric surface roughness, and realistic finite-temperature corrections.

These boundaries delineating the framework are theoretically productive, clearly mapping the threshold between established analytics and requisite future derivations.

\subsection{Outlook and Future Directions}

The formulation of the fractal trace anomaly framework intrinsically divides the observable phenomenology into two distinct technological eras. 

In the immediate term, the electromagnetic consequences of the logarithmically running Casimir coefficient $C(d_s,\ln(d/\ell_*))$---manifesting as the fractal Purcell effect, the anomalous Lamb shift, and the spatially modulated Casimir-Polder potential---are robustly accessible. These CQED probes provide an immediate laboratory pathway to physically map the discrete scale invariance of the vacuum field. Confirming the analytic forms of Eqs.~\eqref{eq:CP_fractal}--\eqref{eq:Lamb_fractal} will firmly establish the effective parameters of the framework without necessitating gravitational measurements.

In the long term, isolating the explicit breaking of conformal symmetry via the semiclassical trace anomaly ($T^\mu{}_{\mu} \neq 0$) requires paradigm-shifting macroscopic engineering. Single-atom gravitational redshifts and quadrupolar splitting remain astronomically suppressed by the Planck area $(l_p/d)^2$. Promoting it to a quantitatively predictive theory necessitates several formal developments. Foremost among these is the exact determination of the crossover coefficient $C_n(d)$ via a full electromagnetic calculation for finite-level self-similar configurations. Furthermore, to move beyond an integrated effective trace, a rigorous local computation of the renormalized stress-energy tensor, $\langle T^\mu{}_{\nu}\rangle$, must be executed for at least one representative self-similar boundary problem.

In parallel, realistic experimental modeling dictates the incorporation of material response and dissipation, which are most naturally addressed by embedding the fractal geometries within a formal scattering approach or a worldline formalism. Finally, a systematic analysis of the asymptotic behavior is required to map the crossover from the discrete prefractal regime to the idealized $n\to\infty$ limit. Only through the completion of these analytical and computational steps can the proposed effective trace be elevated from a plausible ansatz to a rigorous, predictive theoretical framework bridging tabletop quantum optics with the macroscopic generation of semiclassical gravity.

\section{Conclusion}
In summary, a precise theoretical formulation of the ``fractal Casimir effect'' necessitates a strict delineation of the underlying spectral and boundary-value problems.
Previous treatments have frequently conflated three distinct geometric regimes: fields propagating on genuine fractal Laplacians, Euclidean cavities characterized by fractal boundaries, and self-similar arrangements of standard macroscopic plates.
By systematically decoupling these configurations and establishing the appropriate thermodynamic bounds, the present work provides a rigorous foundation for the proposed effective trace framework.

That framework demonstrates the following: if a plate-like self-similar vacuum system is described by a Casimir coefficient that runs logarithmically with scale, then the anisotropic integrated vacuum stress tensor acquires a nonzero trace proportional to $\partial_{\ln d}C$.
When combined with the rigorous thermal trace of fractal radiation, this yields a unified effective trace formula, Eq.~\eqref{eq:total_trace_explicit_new}, and an associated semiclassical curvature source, Eq.~\eqref{eq:R_new}.

This result is not a local anomaly theorem and should not be presented as one.
Its value is instead to provide a sharp, testable organizing principle for future calculations: once a realistic electromagnetic model yields the running coefficient $C_n(d)$, the resulting integrated trace and its possible gravitational interpretation follow immediately.
In that sense, the framework proposed here aims to place the macroscopic backreaction of self-similar vacuum systems on a firmer conceptual footing.



\bibliography{references}
\bibliographystyle{unsrt}

\end{document}
